% GENERATED FILE -- do not edit. The thesis source is the GIU template folder thesis/overleaf/
% (main file bachelor.tex). Run `bash thesis/build.sh` to rebuild this file and thesis/thesis-draft.pdf.
% This single file contains the whole thesis for pasting into Overleaf; upload figures/ for the images.
\documentclass[12pt, a4paper]{book}%, oneside

\usepackage[headings]{fullpage}
\usepackage{setspace}
% ---- Added for this thesis: encoding, maths, long tables, special characters, URLs ----
\usepackage[T1]{fontenc}
\usepackage[utf8]{inputenc}
\usepackage{lmodern}
\usepackage{amsmath,amssymb}
\usepackage{longtable,calc}
\usepackage{etoolbox}
\makeatletter
\patchcmd\longtable{\par}{\if@noskipsec\mbox{}\fi\par}{}{}
\makeatother
\usepackage{footnotehyper}
\makesavenoteenv{longtable}
\usepackage{xcolor}
\usepackage{fvextra}
\fvset{breaklines=true}
\usepackage{newunicodechar}
\newunicodechar{†}{$\dagger$}\newunicodechar{κ}{$\kappa$}\newunicodechar{≥}{$\geq$}\newunicodechar{≤}{$\leq$}\newunicodechar{→}{$\rightarrow$}\newunicodechar{↔}{$\leftrightarrow$}\newunicodechar{±}{$\pm$}\newunicodechar{✔}{Y}\newunicodechar{≠}{$\neq$}
\providecommand{\tightlist}{\setlength{\itemsep}{0pt}\setlength{\parskip}{0pt}}
\setlength{\emergencystretch}{3em}
\renewcommand{\arraystretch}{1.15}
\usepackage{xurl}
% ---- end of additions ----

% \usepackage[top=1.5cm,right=1.5cm,bottom=1.5cm,left=1.5cm]{geometry}
% \usepackage{fancyhdr}
% \setlength{\headheight}{20pt} 
% \usepackage{graphics}
\usepackage{graphicx}
\graphicspath{{figures/}} % added: screenshots live in figures/
% Images: shown when the figures/ folder is next to this file; otherwise a labelled box appears.
\newcommand{\thesisimage}[2][]{\IfFileExists{figures/#2}{\includegraphics[#1]{figures/#2}}{\fbox{\parbox[c][3cm][c]{0.6\textwidth}{\centering\small Image \texttt{\detokenize{#2}} not found: upload the \texttt{figures} folder.}}}}
\usepackage{float}
\usepackage{subfig}
\usepackage{array}
% \usepackage{xtab}
\usepackage{multirow}
\usepackage{booktabs}
\usepackage{appendix}
\usepackage[pdfborder={0 0 0}, colorlinks=true, linkcolor=black, citecolor=blue, urlcolor=blue]{hyperref}
% printonlyused removed: the list of abbreviations is printed in full; nohyperlinks avoids broken links
\usepackage[nohyperlinks]{acronym}
\usepackage{algorithmic}
\usepackage{algorithm}
\usepackage{ifpdf}
% \setlength{\parindent}{0pt}

\newcommand{\submissionDay}{5}
\newcommand{\submissionMonth}{January}
\newcommand{\submissionYear}{2027}
\newcommand{\submissionDate}{\submissionDay~\submissionMonth,~\submissionYear}
\newcommand{\typeOfThesis}{Bachelor Thesis}

\newcommand{\titleOfThesisOne}{JadwAI: AI-Assisted University Timetabling Through Natural-Language Constraint Extraction}

\newcommand{\authorOfThesis}{Rawda Yassin}
\newcommand{\supervisorOne}{Dr. John Zaki}
\newcommand{\supervisorTwo}{}
\newcommand{\supervisorThree}{}

\newcommand{\includefig}[4]{
    \begin{figure}[ht]
     \centering
      \thesisimage[width=#1\textwidth]{#2}
      \caption{#3}
      \label{#4}
    \end{figure}
}

\newcommand{\includefigWSC}[5]{
    \begin{figure}[ht]
     \centering
      \thesisimage[width=#1\textwidth]{#2}
      \caption[#3]{#4}
      \label{#5}
    \end{figure}
}

\newcommand{\includeeps}[4]{
\includefig{#1}{#2.eps}{#3}{#4}
}

\newcommand{\includeepsWSC}[5]{
\includefigWSC{#1}{#2.eps}{#3}{#4}{#5}
}


\ifpdf
\pdfinfo {
	/Author (\authorOfThesis)
	/Title (\titleOfThesisOne)
	/Subject (\typeOfThesis)
	/Keywords ()
	/CreationDate (D:20090707085533)
}
\fi

\begin{document}
% \overfullrule=5pt
\pagestyle{plain}
\pagenumbering{Roman}

\clearpage
% ======== title_page.tex ========
\newcommand{\titlePage}{

\thispagestyle{empty}
\begin{center}
	\thesisimage[width=3.5cm]{giu_logo.png}\\[2mm]
	\textbf{German International University}\\[2mm]
		% NOTE: check the exact faculty name (old draft said "ICS - Software Engineering")
		\textbf{Faculty of Informatics and Computer Science}\\[1mm]

	
	\vspace{2cm}
	\doublespacing
	{\Huge \textbf{\titleOfThesisOne}}\\
	\singlespacing
	\vspace{2cm}
	{\large \textbf{\typeOfThesis}}\\
	
	\vfill
	\parbox{1cm}{
  		\begin{large}
    			\begin{tabbing}
       			Author: \hspace{2cm}  
        			\=\authorOfThesis\\[2mm]
      			Supervisor: 
        			\>\supervisorOne\\[2mm]
      			Submission Date: 
        			\>\submissionDate\\
    			\end{tabbing}
  		\end{large}
	}\\
\end{center}
\clearpage
}
%++++++++++++++++++++++++++++++++++++++++++++++++++++++++++++++++++++
\titlePage
\thispagestyle{empty}\ \clearpage
\titlePage
%++++++++++++++++++++++++++++++++++++++++++++++++++++++++++++++++++++
\thispagestyle{empty}
This is to certify that:
\begin{itemize}
\item[(i)] the thesis comprises only my original work toward the Bachelor Degree
\item[(ii)] due acknowledgement has been made in the text to all other material used
\end{itemize}

\vspace{2cm}
\begin{flushright}
\rule[0mm]{6cm}{0.2mm}\\
\authorOfThesis\\
\submissionDay~\submissionMonth,~\submissionYear\\
\end{flushright}
\clearpage
\clearpage


\clearpage
% ======== acknowledgments.tex ========
\chapter*{Acknowledgments}
\addcontentsline{toc}{chapter}{Acknowledgments}
\label{chap:ack}
% TO WRITE, in your own words: supervisor, anyone who helped, and any funding.
[Acknowledgments to be written.]
\clearpage


\clearpage
% ======== abstract.tex ========
\chapter*{Abstract}
% \addcontentsline{toc}{chapter}{Abstract}
\label{chap:abstract}
% TO WRITE LAST, in your own words, about 250-300 words in three parts (How to Write a Better Thesis, ch. 11):
% (1) why the work was done and what it aimed to achieve;
% (2) what was built and how it was evaluated, with the main results;
% (3) what was concluded.
[Abstract to be written.]
\clearpage

\tableofcontents
% \addcontentsline{toc}{chapter}{Contents}
\clearpage 

\pagestyle{headings}
\pagenumbering{arabic}

% \setlength\parskip{15pt}
\setlength\parskip{.5\baselineskip plus .2\baselineskip
	minus .4\baselineskip}
% \setlength\parskip{.5\baselineskip \@plus .1\baselineskip \@minus ..1\baselineskip}


\clearpage
% ======== sections/introduction.tex ========
% NOTE (migration 23 Nov 2026): Text moved unchanged from main.tex. Much of it was drafted with AI assistance; rewrite it in your own words before submission (the supervisor runs AI checks).
% NOTE (23 Nov 2026): Research questions, scope, contributions and organisation updated to match the rewritten Methodology. RQ3/RQ4 reframing assumes the supervisor accepts RQ-EVIDENCE.md; change it together with Methodology Section 3.1 otherwise.
\hypertarget{introduction}{%
\chapter{Introduction}\label{introduction}}


\hypertarget{problem-statement}{%
\section{Problem statement}\label{problem-statement}}

Administrators must record scheduling requirements accurately enough to protect access to required teaching while accommodating staff availability and suitable facilities. If the system misunderstands a requirement, a timetable can appear acceptable while failing the intended institutional need. The practical problem is therefore how to support clear requirement entry, useful clarification, and informed approval without adding excessive administrative effort.

The five reviewed timetabling papers provide evidence about scheduling needs and capabilities, but do not evaluate ordinary-language requirement entry with this review workflow. This thesis investigates that task within a limited timetabling scope. Its gap claim is bounded by the reviewed sources; a wider documented literature search remains planned.

\hypertarget{research-gap}{%
\section{Research gap}\label{research-gap}}

Within the five retained timetabling papers, constraint acquisition is a developer or
administrator task and is not evaluated as a research problem (Chapter~\ref{literature-review}). The additional LLM-related sources verified so far address
different questions: LLMs solving scheduling problems described in text
\cite{sharma2026schedbench}, a pilot on appliance scheduling in Italian
\cite{sanguinetti2025preferences}, and a benchmark of LLMs writing constraint
models from problem descriptions~\cite{michailidis2025cpbench},
whose task differs because the output is a general constraint program
with no institutional data to resolve against and no review step. The
gap addressed here is therefore the \emph{evaluation of
natural-language-to-constraint extraction for university course
timetabling}, with human review and independent validation, using a
fixed formal model so that extraction quality can be separated from
solver behaviour. The novelty claim is limited to this evaluated
combination; the use of an LLM is not itself claimed to be novel.

\hypertarget{research-questions}{%
\section{Research questions}\label{research-questions}}

\begin{itemize}
\tightlist
\item
  \textbf{RQ1 (primary).} How accurately can an LLM extract formal
  university scheduling constraints from natural-language requirements
  --- including entity references, parameters and hard/soft
  classification --- when output is restricted to a closed constraint
  vocabulary?
\item
  \textbf{RQ2.} Can the system identify ambiguous, missing,
  contradictory or unsupported requirements and request useful
  clarification instead of guessing?
\item
  \textbf{RQ3.} Does AI-assisted constraint entry with human review
  reduce entry time and the number of corrections compared with manual
  form entry while preserving correct interpretation?
\item
  \textbf{RQ4.} How does a locally hosted base language model compare
  with the same model fine-tuned on reviewed, guide-grounded examples,
  in answer quality and latency, on the same held-out cases?
\end{itemize}

RQ1 and RQ2 are evaluated overall and separately for English, formal Arabic, Egyptian Arabic and mixed Arabic/English, by comparing a rule-based parser with a locally hosted language model on development data with draft labels. The results are therefore development-level evidence, not a confirmatory held-out result. RQ3 could not be studied, because no participants or ethics approval were available in the project period; it is discussed as future work. RQ4 originally compared several commercial models; because no budget for paid calls was approved, it compares a local base model with its fine-tuned versions (Chapter~\ref{methodology-proposed-study}).

A supporting engineering question concerns whether the surrounding
system can keep the model honest: \emph{can an independent validator and
an approval workflow ensure that only reviewed constraints reach the
solver and that no schedule with a supported hard-rule violation is
published?} This is verified by testing
(Section~\ref{scheduling-validation-and-system-behaviour-evaluation}), not by statistical inference.

\hypertarget{objectives}{%
\section{Objectives}\label{objectives}}

\begin{enumerate}
\def\labelenumi{\arabic{enumi}.}
\tightlist
\item
  Identify and document the supported educational and administrative scheduling requirements.
\item
  Enable administrators to enter requirements in ordinary language, inspect their interpretation, and clarify or correct them before approval.
\item
  Generate candidate course timetables that respect approved attendance, availability and facility requirements.
\item
  Give administrators understandable checks and feedback for generated timetables and manual changes before publication.
\item
  Evaluate interpretation correctness, administrative task effort and understanding separately from timetable quality.
\item
  Deliver a working prototype, reproducible experiments and this thesis.
\end{enumerate}

\hypertarget{contribution-proportional-to-the-evidence}{%
\section{Contribution (proportional to the
evidence)}\label{contribution-proportional-to-the-evidence}}

The contributions are: (i) a documented set of supported rules and a
machine-readable output contract for timetabling requirements, with
explicit decisions for clarification, unsupported and conflicting
requests; (ii) a synthetic development corpus of 144 statements in four
language categories with a versioned annotation guide, whose labels are
draft and unreviewed; (iii) a development-level comparison of a
rule-based parser and a locally hosted language model on that corpus,
with family-level uncertainty; (iv) a fine-tuning study of a local
assistant model in three rounds, including a round that made it worse
and the reason; and (v) a working prototype for teaching and examination
timetables in which an approval workflow and an independent validator
keep unreviewed or invalid output from being used. No claim is made
about accuracy on real administrators' statements or about reduced
administrative effort.

\hypertarget{scope}{%
\section{Scope}\label{scope}}

\textbf{Confirmed project scope:} JadwAI supports course timetabling for one department and one semester, using a repeating weekly pattern. Lectures, tutorials and laboratories have durations and weekly occurrences; student groups are the unit of student-conflict checking. Natural-language requirements in English, formal Arabic, Egyptian Arabic and mixed Arabic/English are part of the core research scope, alongside an English/Arabic interface with right-to-left presentation. JAI (pronounced ``jay'') is the website's chat assistant; its language model runs on the same server as the application.

The solver is intended to choose instructors, times and rooms. Eligibility is recorded separately for lecture, tutorial and laboratory activities. Each activity retains one instructor across its weekly occurrences; instructor, time and room assignments may be locked. Mandatory weekly instructor workload limits and configurable mandatory or preferred daily limits for instructors and groups are in scope. The confirmed unit is elapsed teaching hours: a 90-minute session contributes 1.5 hours, counted once per instructor irrespective of how many groups attend; fixed breaks are excluded. Configured weekly instructor maxima and global hard daily caps are implemented. Preferred daily-limit variants and entity-specific daily overrides remain proposed extensions.

\textbf{Added during the project:} examination timetabling (quizzes, midterms and finals, with halls and TA invigilators), repair suggestions for impossible timetables (proposed to the administrator, never applied automatically), and generation in parts for university-scale data. \textbf{Deferred extensions:} individual enrolments, minimal-change regeneration, calendar feeds and notifications. \textbf{Excluded:} native mobile applications, attendance, payroll, full student-information-system integration, automated registration or sectioning, and institution-wide production deployment. Approved scope is distinguished from implemented behaviour in Section~\ref{implementation-status}.

\hypertarget{organisation-of-the-thesis}{%
\section{Organisation of the thesis}\label{organisation-of-the-thesis}}

Chapter~\ref{chap:background} introduces university timetabling and the constraint-entry problem. Chapter~\ref{literature-review} reviews five timetabling papers through educational needs, implemented capabilities and reported evidence, and examines a study of natural-language preference interpretation. Chapter~\ref{methodology-proposed-study} describes the method followed: the research design, the interpretation contract, the development corpus, the methods compared, the metrics, the fine-tuning of the assistant model, the engineering verification, and the threats to validity. Chapter~\ref{chap:implementation} describes the implemented system and the JAI assistant. Chapter~\ref{evidence} reports the measurements, and Chapter~\ref{chap:conclusion} summarises the findings, limitations and future work. The appendices list the requirements, screenshots of the prototype and planning records from the proposal.
\clearpage

\clearpage
% ======== sections/background.tex ========
\chapter{Background}\label{chap:background}

% NOTE (migration 23 Nov 2026): Text moved unchanged from main.tex. Much of it was drafted with AI assistance; rewrite it in your own words before submission (the supervisor runs AI checks).
% NOTE (migration 23 Nov 2026): Sections below were moved from the Introduction. Still to add in your own words: constraint programming and CP-SAT; large language models, local hosting (Ollama) and QLoRA fine-tuning; right-to-left and mixed-direction text. Cite sources for each.
\hypertarget{university-timetabling}{%
\section{University timetabling}\label{university-timetabling}}

University timetabling shapes whether students can attend required classes and assessments and whether staff can deliver them within the institution's available time and facilities. An overlapping required lecture can prevent attendance; an unsuitable room can prevent a planned laboratory activity; poorly coordinated examination dates can conflict with local assessment rules. These are educational and administrative problems before they are computational ones. Course timetabling organises repeated teaching activities, while examination timetabling organises assessment events, with different room and calendar requirements~\cite{petrovic2004university}.

The five retained timetabling papers illustrate different priorities. Sobh, Cousens and Patel~\cite{sobh_course_scheduler} focus on offering course combinations that meet students' degree requirements. Müller, Rudová and Müllerová~\cite{muller2025itc} describe the need to attend compatible lectures, seminars and laboratories, including travel between locations. Abdennadher and Edward~\cite{abdennadher_guc_exams} consider examination clashes, retakes, core-exam limits and spacing between difficult examinations at GUC. Kaya~\cite{kaya2026drl} includes institutional-scale examination revision and coordinator control.

A mandatory rule protects a requirement that the institution has decided must hold, such as avoiding simultaneous classes for the same student group. A preference expresses a desirable arrangement, such as a preferred teaching period, which may compete with other needs. These categories and their priorities must be agreed locally. The papers primarily evaluate scheduling results; they do not justify assuming that a conflict-free timetable automatically improves learning or graduation. This study accordingly treats educational access and administrative usability as motivations and measures the outcomes within its actual scope.

\hypertarget{the-constraint-entry-problem}{%
\section{The constraint-entry
problem}\label{the-constraint-entry-problem}}

Solvers need explicit rules: entities, parameters, severities and
weights. Administrators, by contrast, express requirements informally
and often in several rules at once. A statement such as "Dr Ahmed cannot
teach on Tuesdays and prefers morning classes" contains a mandatory
unavailability and a preference; "morning" has no meaning to a solver
until an institution defines it. A wrong reading --- treating the
preference as mandatory, resolving the name to the wrong person, or
choosing an arbitrary definition of "morning" --- can produce a
timetable that satisfies the formal model while failing the
administrator\textquotesingle s intention.

In the five timetabling papers reviewed, this translation step is not itself
studied. Constraints are encoded by developers in declarative programs
\cite{abdennadher_guc_exams},
supplied in standard benchmark formats~\cite{muller2025itc}, formalised
by researchers for a given institution~\cite{kaya2026drl}, set as weights by
the timetable officer~\cite{petrovic2004university}, or entered by
instructors choosing time slots~\cite{sobh_course_scheduler}.
Existing systems offer form-based entry, which is precise but laborious.
An observation recorded in the project specification from a walkthrough
of the public UniTime demonstration on 4 October 2026 (a sample dataset
and a single session, not a systematic evaluation, and not independently
re-checked here) illustrates the effort: recording that one instructor
prefers no teaching after 3 PM required navigating to the instructor,
opening preferences, choosing a severity and marking 15 half-hour cells
across five weekdays. This observation motivates the project but is not
evidence about any deployed system.

\hypertarget{motivation-for-ai-assisted-extraction}{%
\section{Motivation for AI-assisted
extraction}\label{motivation-for-ai-assisted-extraction}}

Large language models can read free text and, when constrained to a
schema, produce structured output. This suggests a workflow in which an
administrator writes a requirement, a model proposes formal rules, and
the administrator reviews them. The suggestion carries risks that the
design must address rather than assume away: schema-valid output is not
necessarily a correct interpretation; models may invent entities or
silently choose thresholds; and a correct extraction does not make a
schedule valid. Recent work on LLMs in scheduling indicates that wording
affects constraint faithfulness when models are asked to \emph{produce}
schedules from text~\cite{sharma2026schedbench} (preprint), and that
small-scale attempts to convert user preferences into optimisation
constraints in other domains are still preliminary~\cite{sanguinetti2025preferences}
(preprint). These sources concern different tasks
and settings; they motivate care and measurement, not a conclusion about
university timetabling.
\clearpage


% add more chapters here
\clearpage
% ======== sections/related_work.tex ========
% NOTE (migration 23 Nov 2026): Text moved unchanged from main.tex. Much of it was drafted with AI assistance; rewrite it in your own words before submission (the supervisor runs AI checks).
\chapter{Related Work}\label{literature-review}

\section{Purpose and review approach}
This chapter examines university timetabling from an educational and administrative perspective. The main questions are what problem each paper addresses, who benefits, what capabilities were implemented or studied, which educational requirements were considered, what the evaluation establishes, and what remains missing. Algorithms are mentioned only where necessary to understand a capability or a limitation.

The review is a focused reading of five retained timetabling papers and one additional paper on natural-language preference interpretation, not a systematic review of all timetabling research. Two concern examination timetabling, one concerns a course-scheduling tool, one analyses real-world course-timetabling benchmarks, and one surveys university timetabling. These are different kinds of evidence: a review chapter does not demonstrate a single deployed product, a benchmark study does not constitute a user study, and a prototype demonstration does not establish improved academic outcomes. The discussion therefore distinguishes implemented functionality, author-reported scheduling results, intended educational benefits, and future work. Page references below refer to the supplied PDF page order where journal pagination is unavailable. Incomplete publication metadata is retained as provisional rather than guessed.

\section{Course planning and access to required teaching}
\subsection{Sobh, Cousens and Patel: Course Scheduler}
\textbf{Educational problem and beneficiaries.} Sobh, Cousens and Patel~\cite{sobh_course_scheduler} address the difficulty of deciding which courses to offer and when to offer them so that students can meet degree requirements. Their motivating example is a senior student with only a small set of remaining courses: an unsuitable offering or a clash can obstruct completion of the degree. The primary operator is the department chairperson or administrator; instructors contribute availability and teaching preferences, while students benefit indirectly through access to suitable course combinations (supplied PDF, pp.~1--2).

\textbf{Implemented capabilities and educational requirements.} The described suite combines student course requirements with instructor teaching capabilities and load limits to identify possible course offerings. Instructors select teaching times and rank possible course sets. Administrators inspect schedule viability, view specific student conflicts, and adjust course times through the interfaces. The tool gives feedback on which student requirements a proposed schedule satisfies. Thus, its educational emphasis is course access and progression, rather than merely filling a weekly grid (pp.~12--17).

\textbf{Evidence and limits.} The paper illustrates the interfaces and shows example combinations with 100\% and 82\% viability, including a conflict explanation for an individual student. These are demonstration results, not a controlled evaluation of graduation rates, student satisfaction, or staff time savings. Helping students graduate sooner is a motivation, not a measured outcome. The authors explicitly state that the suite does not address classroom allocation or class-size issues. They also identify fixed time slots and dependence on instructors' chosen times as limitations; generating times from student needs is future work (pp.~16--17).

\textbf{Implications for this project.} A useful timetable must protect access to required teaching and explain conflicts in terms an administrator can act on. This supports conflict reporting, review, and safe manual adjustment in the proposed project. However, the proposed prototype starts with predefined course activities and student groups. It does not select semester course offerings, optimise degree completion, or provide individual enrolment-based schedules unless those extensions are approved.

\subsection{Müller, Rudová and Müllerová: real-world course timetabling}
\textbf{Educational problem and beneficiaries.} Müller et al.~\cite{muller2025itc} examine the diversity of real university course-timetabling needs through ITC~2019. The central educational issue is whether students can attend the required combination of lectures, seminars, and laboratories. Beneficiaries include students seeking compatible class combinations and administrators coordinating teaching across faculties, buildings, and institutional calendars.

\textbf{Capabilities represented in the study.} The paper describes thirty benchmark instances from ten universities using UniTime in practice. The competition model combines class time and room assignment with student sectioning: students must be assigned the appropriate class types and compatible lecture--seminar or lecture--laboratory combinations. It represents room capacities and unavailable periods, repeated meetings and week patterns, travel time between rooms, and requirements on gaps, breaks, teaching days, and daily load (journal pp.~248--250). A student conflict can arise from either overlapping classes or insufficient travel time. These are requirements represented by the benchmark and contributing institutions; the paper is not presenting a newly tested interface for all these functions.

\textbf{Educational significance.} The institution descriptions show why a single definition of timetable quality is inadequate. Examples include travel between sports facilities, joint programmes creating unusual course combinations, and concentrated attendance days for students combining study with work. Scheduling therefore supports participation in a programme under practical attendance conditions, while institutional priorities determine acceptable trade-offs.

\textbf{Evidence and limits.} The paper analyses instance characteristics and compares scheduling solutions. This provides substantial evidence that real timetabling requires more than avoiding simultaneous room use. Its evaluation concerns feasibility and scheduling penalties, including student conflicts; it does not establish improvements in learning, retention, or satisfaction. Competition results also cannot be transferred directly to a smaller prototype with a different scope.

\textbf{Implications for this project.} The proposed system should make its supported educational requirements explicit and assess them against the approved institutional policy. Initial group-based conflict checking is a deliberate simplification of individual sectioning. Travel-time checking, irregular week patterns, and student sectioning should not be claimed as implemented features of the prototype without corresponding requirements and tests.

\section[Examination access and workload]{Examination access, workload and institutional coordination}
\subsection{Abdennadher and Edward: examination timetabling at GUC}
\textbf{Educational problem and beneficiaries.} Abdennadher and Edward~\cite{abdennadher_guc_exams} describe examination scheduling at the German University in Cairo. At the time described, manual course scheduling reportedly took about three weeks and examination scheduling about one week. Shared courses and resources across faculties made separate faculty schedules unsuitable. Examination administrators are the direct users; students benefit from schedules that allow them to sit their examinations, including examinations retaken from earlier semesters (supplied PDF, pp.~1--3).

\textbf{Implemented capabilities and educational requirements.} The system assigns examinations to periods, distributes students among rooms and seats, and displays results through a graphical interface. Its requirements include avoiding simultaneous exams for the same group, preventing semester clashes within a major to accommodate retakes, allowing no more than one core exam per day, limiting daily examinations, separating difficult examinations, respecting available seats, and supporting fixed and unavailable periods. Difficulty ratings are described as supplied by lecturers and students. These features connect scheduling to assessment workload and access; they should not be interpreted as evidence that the software itself determines pedagogically appropriate difficulty ratings (pp.~3--7).

\textbf{Evidence and limits.} The authors report a test using Spring~2007 final-examination data: 140 courses, 5,500 students and 1,300 available seats, with a solution generated within a few seconds. This demonstrates scheduling functionality on the stated data. The historical manual preparation time is not a measured before-and-after comparison of the entire administrative workflow. The paper does not report effects on exam performance, stress, or satisfaction. Integration with the invigilation system and further timetable repair and optimisation are future work (pp.~7--8).

\textbf{Implications for this project.} The paper supports clear distinctions between student access rules, assessment-spacing rules, and resource restrictions. It also motivates collecting local institutional requirements rather than copying another university's rules. Examination spacing and retake accommodation are relevant if an examination extension is approved; they are not automatically requirements for the current weekly course-timetabling scope.

\subsection{Kaya: multi-campus examination decision support}
\textbf{Educational problem and beneficiaries.} Kaya~\cite{kaya2026drl} addresses large-scale examination scheduling at a multi-campus Turkish university. Examinations have different durations, and many cohorts require more seats than a single room provides. Coordinators need to organise examinations within campus and facility restrictions while maintaining compatible assessment schedules. Students and teaching staff are affected by the resulting placements, although coordinators are the direct decision-support users.

\textbf{Implemented capabilities and educational requirements.} The reported framework handles duration-based placements, allocation across multiple rooms, room types and capacities, campus restrictions, instructor availability, and cohort non-overlap. Its institutional decision-support workflow lets coordinators reopen selected examinations and revise placements while preserving locked assignments. The paper describes selecting courses, moving placements, accepting suggestions, and approving a schedule (Sections~4.7 and 5.11--5.12). These capabilities matter because institutional timetables are revised and reviewed rather than simply generated once.

\textbf{Evidence and limits.} The main comparison uses 1,127 representative examinations, 110 rooms and 520 instructors, derived from 1,332 raw events. The paper reports high assignment coverage and an independent audit finding no hard-rule violation among assigned examinations. Coverage and validity must be read together: unassigned exams still require resolution. Crucially, cohort conflicts were made active through synthetic conflict injection because naturally occurring cohort overlap was absent from the evaluated production period; an accessibility requirement was also inactive. The operational logs establish execution at institutional scale, but the forty recorded coordinator interventions were a controlled replay, not observed live coordinator behaviour. The study is limited to one institution and period, and it does not establish reduced student stress, improved marks, or live administrative time savings (Sections~5.12 and 6.3).

\textbf{Implications for this project.} This paper supports reviewed changes, preservation of approved decisions, explicit reporting of unscheduled activities, and independent checking after edits. It also demonstrates why an educational review must ask whether relevant requirements were exercised in realistic data. The proposed evaluation will distinguish operational records, simulated cases, and participant evidence rather than treating them as interchangeable.

\section{University timetabling as an institutional decision}
\subsection{Petrovic and Burke: University Timetabling}
\textbf{Problem and beneficiaries.} Petrovic and Burke~\cite{petrovic2004university} review university course and examination timetabling as a major administrative activity. Their chapter concerns the interests of timetable officers, students and teachers, particularly the difficulty of satisfying different institutional priorities within limited resources (supplied PDF, pp.~2--3).

\textbf{Capabilities discussed.} The chapter reviews approaches to constructing feasible schedules, balancing several quality criteria, reusing previous scheduling experience, and adapting methods across problems. Its educational contribution to this review is the distinction between mandatory attendance and resource rules and desirable preferences, alongside recognition that those rules vary across universities. It explains why course and examination scheduling need different assumptions about repeated meetings and room use. Because this is a review chapter, these capabilities should not be presented as features implemented together in a single system.

\textbf{Evidence and limits.} The reviewed studies provide scheduling results and examples of ways to represent competing objectives. The chapter does not supply one controlled evaluation of an integrated product or evidence of improved learning outcomes. Its range of examples supports context-sensitive requirements, not a universal checklist that every university should adopt.

\textbf{Implications for this project.} The proposed prototype should let administrators express supported local requirements, identify which are mandatory, and inspect which preferences remain unmet. The distinction is important educationally: avoiding a clash protects access, while satisfying a preferred teaching time improves convenience under a particular institutional policy. Neither should be conflated with a demonstrated change in educational achievement.

\section{Natural-language preferences and faithful interpretation}
\subsection{Sanguinetti et al.: From User Preferences to Optimization Constraints Using Large Language Models}
\textbf{Problem and beneficiaries.} Sanguinetti et al.~\cite{sanguinetti2025preferences} investigate how people can express household appliance needs in ordinary language so that those needs become constraints for an energy optimizer. The motivation is that a community's energy-saving objectives may conflict with individual routines: a technically efficient recommendation is of limited use if it ignores when a person needs hot water or air conditioning. The direct beneficiaries are household users in Italian renewable energy communities. Although this is not an educational study, the same concern motivates the proposed timetable interface: an administrator's intended attendance or availability requirement must survive its translation into a formal rule (supplied PDF, pp.~1--3).

\textbf{Capabilities studied.} The evaluated component converts Italian utterances into representations of appliance operating times or desired temperatures. It supports requests applying throughout the time horizon, within a specified interval, after a start time, or until an end time. The input includes both the utterance and manually identified preference spans marked with XML tags. The study therefore evaluates conversion of identified preferences, rather than the complete discovery of all requirements in unrestricted text. A connection to an external optimizer is described within the wider project; the paper does not establish a completed, evaluated workflow with user clarification, interpretation approval and downstream schedule validation (pp.~3--8).

\textbf{Evidence.} The authors manually annotate a pilot subset of 26 utterances drawn from a 157-utterance Italian corpus. Prompts provide no example, one example, or five examples. Evaluation separates character-level similarity (ChrF) from correctness of variables and time conditions. In the five-example setting, Table~1 reports variable accuracy of 0.8571 and condition accuracy of 0.4286 for LLaMAntino-3-ANITA; the corresponding ChatGPT values are 0.7619 and 0.3571. These are the paper's component-level measures, not the proportion of complete requests interpreted correctly. Conditions remain substantially harder than variable identification, and adding examples does not improve every component consistently. The authors conclude that the tested prompting approach is inadequate for this task and propose further data collection and supervised fine-tuning (pp.~5--10).

\textbf{Limits of the evidence.} The small, handcrafted sample, Italian appliance domain and supplied preference spans limit transfer to English university requirements containing multiple rules, ambiguous staff names or institutional terms. Character similarity cannot establish that a constraint preserves user intent. The supplied version also has a reporting inconsistency: its model-selection and results narrative describes four models and excludes two, while Table~1 additionally lists Cerbero, Maestrale and IT5 without matching experimental descriptions. Those additional rows are not used here to infer a reproducible ranking. The paper does not report a controlled comparison of manual and assisted entry, measured user effort, educational outcomes, or an evaluated clarification-and-approval process. Functional evaluation after connection to the optimizer is proposed as future work (pp.~7--10).

\textbf{Implications for this project.} This is a close conceptual precedent for converting everyday requirements into a restricted constraint vocabulary, but it does not establish that the proposed university application will work. The thesis will consequently assess complete semantic interpretation and individual fields separately, including time boundaries, entity references and mandatory versus preferred status. Requirements will be evaluated from the full statement, with the supplied institutional context; manually highlighted spans will not be assumed in the main evaluation. Review and targeted clarification are proposed safeguards whose effectiveness must be measured. The human comparison will assess final interpretation correctness alongside time and corrections, rather than treating a fluent response or a valid timetable as evidence that the administrator's intent was understood. These implications support the existing SRS requirements for entity resolution, named periods, clarification, review and evaluation (FR-016--FR-020, FR-049 and FR-050); they add no appliance-related functions to the project.

\section[Cross-paper synthesis]{Cross-paper synthesis and implications for the proposed system}
\textbf{Access and progression.} Sobh et al. emphasise course requirements and degree progression, Müller et al. emphasise attendance across compatible class combinations, and the GUC study explicitly considers retake students. Together, they show that avoiding clashes has educational meaning: students must be able to attend required teaching or assessment. However, offering selection, individual sectioning and group-based conflict prevention are different capabilities. The proposed prototype initially addresses the last of these.

\textbf{Workload and assessment organisation.} GUC's core-exam and difficulty-spacing requirements show how assessment policy affects scheduling. These rules are institution-specific; they do not establish an optimal universal spacing policy. Course gaps, late sessions and daily-load balance are related concerns for the proposed course prototype, but their thresholds and priority require local agreement.

\textbf{Staff coordination and usable decisions.} Course Scheduler's conflict feedback and Kaya's reviewed revisions show the importance of giving administrators actionable information and control. A useful system should expose incomplete schedules, violated preferences and the consequences of a change. Faster generation alone is insufficient if the administrator cannot understand or approve the result.

\textbf{Evidence quality.} Across the papers, most evaluation concerns scheduling feasibility, coverage, runtime or penalties. Those outcomes are relevant to institutional operations, but they do not directly measure learning, retention, graduation, fairness, stress or satisfaction. The thesis will therefore describe improved access and workload organisation as intended educational benefits, and claim improvements only for outcomes actually measured.

\textbf{Proposed functional response.} The project will support entry of institutional requirements in ordinary language, display the proposed interpretation, request clarification where needed, and require administrator approval. It will generate course schedules from approved rules; identify group, instructor and room conflicts; check capacity and facility suitability; support comparison and reviewed manual changes; and release validated timetables to authorised users. These are proposed capabilities linked to the SRS, not claims that a completed application already exists. Exam-specific and individual-enrolment features remain extensions.

\section{Research gap and limits of the review}
Within the five retained timetabling papers, the educational needs are represented as predefined rules, input data, preferences or choices made through interfaces. None evaluates the proposed combination of ordinary-language requirement entry, explicit interpretation review, useful clarification, and separate measurement of interpretation accuracy for university course scheduling. This bounded observation motivates studying whether administrators can record their scheduling needs with less effort while preserving their intended meaning. It is not a claim that no related system exists anywhere. Within the sources examined, Arabic and Egyptian Arabic administrative constraint extraction is not evaluated; the Cairo examination context does not itself establish natural-language input support. This motivates language-stratified evaluation without asserting broader novelty.

The additional study by Sanguinetti et al.~\cite{sanguinetti2025preferences} directly evaluates preference-to-constraint conversion in another domain, so the translation task itself is not claimed as new. It does not evaluate the proposed university review workflow. Other sources cited in the introduction provide further context rather than evidence of educational effectiveness in this project. A documented search of further educational timetabling and human-centred scheduling literature is needed before making a broader novelty claim. The search should record sources, dates, inclusion criteria and screening decisions, and include terms concerning student access, administrative workload, participatory scheduling and requirement entry alongside natural-language scheduling.

\section{Summary}
The reviewed literature identifies three priorities for this study: protecting access to required teaching, expressing institution-specific scheduling needs, and keeping administrators informed and in control. Existing evidence mainly establishes scheduling capabilities and operational results. The methodology therefore evaluates both the correctness of recorded requirements and the effort and understanding involved in the administrative workflow, without claiming unmeasured improvements in student learning.
\clearpage

\clearpage
% ======== sections/methodology.tex ========
% NOTE (23 Nov 2026): Rewritten as the method actually followed. Drafted with AI assistance at the owner's request; reword in your own words before submission.
% NOTE (23 Nov 2026): Assumes the supervisor accepts the reframing in RQ-EVIDENCE.md (RQ1 as a development-level comparison, RQ3 as future work, RQ4 as base versus fine-tuned local model). Change Section 3.1 and the Introduction together if the supervisor decides otherwise.
\chapter{Methodology}\label{methodology-proposed-study}

This chapter describes how JadwAI was built and evaluated. The project followed a design-and-evaluate approach: a working prototype was implemented and then measured on synthetic data. Section~\ref{meth-design} states how each research question was answered and how this differs from the original proposal. The later sections describe the interpretation contract, the development corpus, the extraction methods compared, the scoring, the JAI assistant and its fine-tuning, the engineering verification of the scheduler and validator, and the records kept for reproducibility. The measurements themselves are reported in Chapter~\ref{evidence}.

\section{Research design}\label{meth-design}

The study has three parts, kept separate so that one kind of evidence is not read as another.

\begin{enumerate}
\item \textbf{Offline interpretation experiment} (RQ1, RQ2). Two interpretation methods, a rule-based parser and a locally hosted language model, are run on the same fixed set of statements in four language categories. Their outputs are scored against draft expected answers by the same scorer.
\item \textbf{Assistant and fine-tuning evaluation} (RQ4). The base local model and three fine-tuned versions of it answer the same held-out assistant cases. They are compared with automatic checks, latency measurements and a blind side-by-side rating.
\item \textbf{Engineering verification} (supporting question). The solver, the independent validator, the repair diagnosis and the approval and publication gates are tested with deterministic test oracles and measured on synthetic and adjusted university-scale data.
\end{enumerate}

A valid timetable is not evidence of a correct interpretation, and a correct interpretation does not guarantee a good timetable. For this reason interpretation, scheduling and validation are measured separately.

\paragraph{Changes from the proposal.} The proposal planned a confirmatory comparison of several commercial language models on a sealed, independently annotated held-out set, and a participant study comparing manual and AI-assisted entry. Neither was possible within the project. No budget was approved for paid model calls, no second annotator was available, the expected answers were not reviewed by a human, and no ethics approval or participants were obtained. The research questions were therefore answered as follows (Table~\ref{tab:rq-method}).

\begin{table}[h]
\centering
\small
\begin{tabular}{@{}p{0.2\textwidth}p{0.36\textwidth}p{0.36\textwidth}@{}}
\toprule
Question & Method used & Change from the proposal \\
\midrule
RQ1 extraction accuracy & Rule-based parser and local Qwen2.5-7B on 144 development statements, same scorer & Development-level comparison on draft labels, not a confirmatory held-out result \\
RQ2 clarification and unsupported requests & Clarification and unsupported-request scores from the same runs; clarifying behaviour of JAI on held-out cases & Exploratory; small samples \\
RQ3 effort compared with manual entry & Not conducted & Moved to future work (Section~\ref{meth-rq3}) \\
RQ4 model comparison & Base local model against the same model fine-tuned in three rounds, on identical held-out cases & Local base versus fine-tuned model instead of several paid providers; cost per call is zero \\
Engineering question & Tests, seeded faults and measured runs & As planned \\
\bottomrule
\end{tabular}
\caption{How each research question was addressed.}
\label{tab:rq-method}
\end{table}

\paragraph{What the study does not claim.} It does not compare JadwAI's solver with other timetabling systems, it does not claim that the prototype covers what systems such as UniTime cover~\cite{muller2025itc}, and it does not estimate how well any method would perform on unseen statements written by real administrators.

\section{The system under study}\label{meth-system}

The primary user is the timetable administrator. Instructors and students only view their own published timetable. The administrator enters academic data (instructors, rooms, student groups, courses and their lectures, tutorials and laboratories), then enters rules either through forms or in plain English, formal Arabic or Egyptian Arabic. The workflow that the evaluation relies on is:

\begin{enumerate}
\item A typed statement is converted into a \emph{proposal}: zero or more structured rules, plus any questions that must be answered and any parts that cannot be represented.
\item The administrator reviews, corrects or rejects the proposal. Only approved rules are stored and used. A proposal from either the parser or a language model is never applied automatically.
\item The solver (Google OR-Tools CP-SAT~\cite{google_cpsat}) generates a timetable from the approved inputs.
\item An independent validator, written separately from the solver model, recomputes every supported rule on the saved timetable. A timetable is published only if it passes validation against the current approved inputs.
\end{enumerate}

This separation is what makes it possible to use an imperfect interpretation method safely. An interpretation error can lead to a wrong rule being proposed. It cannot lead to an unreviewed rule being used or to a timetable that breaks an approved hard rule being published. Implementation details are given in Chapter~\ref{chap:implementation}.

\section{Supported rules and the interpretation contract}\label{formal-representation-of-supported-constraints}

\paragraph{Supported rules.} The interpretation methods may propose only the following:
\begin{itemize}
\item instructor unavailable or preferred times;
\item room unavailable times;
\item group unavailable times;
\item permitted times for an activity or a group;
\item a weekly teaching-hour limit for an instructor;
\item a change to the global daily hour limits for students or instructors;
\item a lock on an existing assignment.
\end{itemize}
Each rule is either mandatory or preferred. Times are stored in whole or half hours. Structural requirements, such as complete assignment, no room, instructor or group overlaps, eligible instructors and suitable rooms, always apply, whichever way the rules were entered.

\paragraph{Output contract.} Every interpretation is a JSON object with four parts: a \emph{decision}, a list of \emph{proposals}, a list of \emph{clarification codes} and a list of \emph{unsupported codes}. The decision is one of five values:
\begin{itemize}
\item \texttt{ready}: every part of the statement is represented by supported proposals;
\item \texttt{clarify}: a day, time, person, threshold or meaning is missing or ambiguous, and a question must be asked instead of guessing;
\item \texttt{unsupported}: no supported rule can represent the request faithfully;
\item \texttt{conflict}: the request is understood, but its hard conditions cannot all hold;
\item \texttt{mixed}: some parts are supported and others are unresolved or unsupported.
\end{itemize}
Clarification and unsupported codes are short fixed labels, for example \texttt{day\_required}, \texttt{undefined\_early} or \texttt{ambiguous\_instructor}. They identify the question, not its exact wording.

\paragraph{Correctness rule.} A proposal is correct when its rule type, target entity, mandatory or preferred strength, day and times (or numeric value) are all equal to the expected proposal. Rule titles are free text and are not scored, and the order of proposals is ignored. A duplicate proposal counts as an extra prediction. A case is exactly correct only when the decision, the complete set of proposals and both code lists all match.

\section{Development corpus}\label{dataset-construction}

\paragraph{Releases.} Two releases of a synthetic corpus exist.
\begin{itemize}
\item Release \texttt{0.1.0-draft} (25 October 2026) contains 128 statements in 32 intent families, with 96 development and 32 reserved statements.
\item Development bundle \texttt{0.2.0-draft} (1 November 2026) keeps the 96 development statements unchanged and adds 48 statements in 12 new families for a second fictional institution, a design school. The reserved statements were not read, moved or relabelled.
\end{itemize}
All interpretation results in this thesis use the 0.2.0 development bundle.

\paragraph{Content.} The bundle has 144 statements in 36 families. Each family is one intent written in four language categories: English, formal Arabic, Egyptian Arabic and mixed Arabic/English, giving 36 statements per category. The first context (96 statements) is a fictional department with a Sunday--Thursday week. The second (48 statements) is a design school with a Saturday--Thursday week and a 12-hour regular teaching load with overtime allowed. Each context lists instructors with aliases, rooms, groups, activities, the named teaching periods and a fictional saved timetable for lock requests.

The expected decisions are \texttt{ready} for 76 statements, \texttt{clarify} for 44, \texttt{unsupported} for 16 and \texttt{conflict} for 8. The statements cover:
\begin{itemize}
\item clear and compound requests;
\item negation, including Egyptian Arabic forms;
\item missing days or times;
\item ambiguous and unknown names, for example two instructors called Nour, and a person absent from the context;
\item undefined terms such as ``early'';
\item hard contradictions;
\item requests outside the supported rules;
\item minute-precise times;
\item conversion from sessions to teaching hours.
\end{itemize}

\paragraph{Labels.} The statements and expected answers were written with the help of an AI assistant, following annotation guide~v2. The guide defines each decision, the supported proposal boundary and how ambiguity is to be resolved. Every item records its origin and has the status \texttt{draft\_unreviewed}. No human annotation, second annotator or adjudication took place, although the tools for them were built (gold-free review exports, immutable review records, Cohen's $\kappa$ comparison and adjudication forms). The expected answers are therefore one author's interpretation, and results against them are development figures.

\paragraph{Reserved material is not a blind test.} Reserved statements are excluded from default tooling, and access requires a logged reason. However, their author has seen the labels, so they cannot serve as an independent held-out set. A fresh, independently written and annotated set would be needed for a confirmatory study. Families, not single statements, are the unit for splitting and for confidence intervals, because four translations of one intent are not independent observations.

\section{Interpretation methods compared}\label{baselines}

Both methods receive exactly the same input for each statement: the statement text, the context's entity records (identifier, kind, English name and code) and the calendar. Arabic aliases from the context are not supplied, so Arabic names must be matched to English records by the method itself.

\paragraph{Rule-based parser.} JadwAI's built-in parser uses deterministic patterns for English, formal Arabic and Egyptian Arabic, and matches names across scripts by their consonant sounds. It was measured twice. After the first run, its clarification codes were aligned with guide~v2, three cases where it guessed instead of asking were changed, and missing negations were added (Section~\ref{evidence-baseline}). Because these changes were made after looking at development errors, its reported figures are optimistic for this corpus.

\paragraph{Local language model.} The language-model method uses \texttt{qwen2.5:7b-instruct} (Qwen2.5-7B-Instruct, 4-bit quantised) served by Ollama on the same laptop, with no network access needed. The statement is sent through the same code path as the running application (\texttt{extract\_candidate} with the local provider). The prompt carries the same task instructions as the adapter for an external model. Decoding is constrained by a JSON schema, so the reply always has the contract's shape. Two settings of this schema matter:

\begin{itemize}
\item \textbf{Field order.} The model generates fields in the order the schema lists them. An early run that asked for the decision first produced \texttt{unsupported} for almost every statement. The schema now lists proposals first, then clarifications and unsupported codes, and the decision last, so the decision is written after the model has stated what it found.
\item \textbf{Code format.} Codes must be short lower-case identifiers of 3 to 60 letters, digits and underscores, starting with a letter. This prevents free-text explanations inside the code lists.
\end{itemize}

On a small check of 8 statements, the reordering raised correct decisions from 1 to 5. These 8 statements are part of the development corpus, so the change was also tuned on development data. Decoding used temperature~0, seed~42 and an 8{,}192-token context window. Each statement had one attempt with a 600-second time limit; a failed or invalid reply counts as an error, not as a skipped case.

\paragraph{Shared format adapter.} Before scoring, both methods' outputs pass through the same adapter. It removes empty optional fields and writes whole hours as integers (10 rather than 10.0), which the contract requires. It does not change half-hour values or any meaning.

\paragraph{External model adapter.} An adapter for a commercial model (GPT-5.4 mini) was implemented and tested with recorded responses, but it was never called because no budget was approved. Therefore no commercial model is compared.

\section{Scoring and metrics}\label{metrics}

All outputs are scored by the project's offline scorer, which first validates each output against the contract and the context. An output that names an entity not in the context, or breaks the contract, counts as wrong. Missing outputs fail their case, and every case stays in the denominator. The following measures are reported overall and separately for each language category:

\begin{itemize}
\item \textbf{Exact output} (primary): the proportion of statements whose whole output matches the expected one.
\item \textbf{Decision accuracy}: the proportion with the correct decision.
\item \textbf{Proposal precision, recall and F1}: proposals matched one-to-one with expected proposals.
\item \textbf{Field F1}: agreement on individual fields of matched proposals, as a looser diagnostic.
\item \textbf{Clarification F1} and \textbf{unsupported-request F1}: agreement on the code lists (RQ2).
\item \textbf{Latency}: wall-clock seconds per statement, reported as the median.
\end{itemize}

Confidence intervals are 95\% percentile bootstrap intervals with 1{,}000 replicates, resampling whole families (seed~42), so that the four translations of an intent stay together. With 36 families the intervals are wide, and the per-language figures, which rest on 36 statements each, are exploratory. No significance test between the two methods is used as a main result, because the labels are draft and the parser was adjusted on the same data.

\section{The JAI assistant and its fine-tuning}\label{meth-jai}

RQ4 was reframed as a comparison between a locally hosted base model and the same model after fine-tuning. The task is JAI's conversational role: explaining JadwAI, answering from the timetable facts supplied in the prompt, asking one focused question when a request is vague, declining off-topic requests, preparing at most one action, and replying in the user's language.

\paragraph{Training data.} Examples are fictional conversations in English, formal Arabic and Egyptian Arabic, grouped into families (one intent with several wordings). Whole families are assigned either to training or to the held-out set, so that no wording of a held-out family is seen in training. The prompts are built by the same code that builds the live chat prompt, including the matching entries of a verified guide (Section~\ref{evidence-jai}). From round~3 on, the dataset builder refuses any example whose answer states a fact about JadwAI that is not in that example's guide entries. The examples were drafted with an AI assistant and approved in bulk by the project owner after reading a change list, not item by item.

\paragraph{Training.} Each round trained a QLoRA adapter (rank~16) on Qwen2.5-7B-Instruct with the Unsloth library on a free T4 GPU (Google Colab or Kaggle), with the loss computed only on JAI's final answer. Each adapter was merged into Ollama's 4-bit weights. To stop the small weight changes being lost when the weights were quantised again, each tensor was re-quantised to a different format (4-bit blocks to 6-bit, 6-bit to 5-bit). A separate script compared the merged weights with the base and the adapter, to check how much of the trained change was present in every trained matrix. Three rounds were trained (Table~\ref{tab:jai-rounds}): v1 on the original 210 examples, v2 on the same answers with the guide added to the prompts and an extra epoch, and v3 on 243 examples whose answers were rewritten to follow the guide. Round~4 data (317 training and 85 held-out rows) was prepared from the failures found in live use but not trained in time.

\paragraph{Evaluation.} All candidates answered the same held-out cases with the production prompt, at temperature~0 and seed~42, on the same laptop. The cases were:
\begin{itemize}
\item 16 conversation cases and 5 action-choice cases from held-out families;
\item for v1 and v3, three wordings of a new question that should be answered with ``not sure''.
\end{itemize}
Three kinds of measurement were used:
\begin{enumerate}
\item \textbf{Automatic checks}: whether the reply is valid, whether it proposes an action it should not, whether it is in the user's language, whether it is at most 80 words, and whether the action chosen in the action cases is correct. These checks do not judge whether an answer is true.
\item \textbf{Latency}: mean seconds per answer.
\item \textbf{Blind side-by-side rating} of v1 and v3 on the 19 conversation answers. The two answers to each case were shown in random order without model names, and each was scored from 1 to 5 on four criteria: correct, logical, helpful and natural, giving at most 20 per case. The key was kept in a separate file. The rater was the project owner, who had read a short summary of some replies beforehand, so the rating is neither independent nor fully blind.
\end{enumerate}
Finally, the preferred candidate (v3) was used in the running application with questions taken from its weakest cases, to see whether the held-out results held in real multi-turn use.

\section{Scheduling and validation evaluation}\label{scheduling-validation-and-system-behaviour-evaluation}

\paragraph{Instances.} Three kinds of data were used:
\begin{itemize}
\item small hand-built instances, each designed to exercise one rule;
\item a demonstration department workspace of 448 weekly sessions, 93 groups, 123 instructors and 65 rooms, derived from real course structures with invented names;
\item the Winter 2026 university workbook (2{,}711 weekly sessions), entered through the data template.
\end{itemize}
The university data needed documented adjustments before a timetable could exist, and is therefore used to test the mechanism of generating in parts, not to produce a real university timetable.

\paragraph{Settings and measures.} All runs used one laptop (Intel Core i5-1038NG7, 16\,GB memory, CPU only), CP-SAT with eight search workers and seed~42. The measures recorded were:
\begin{itemize}
\item whether every session was placed;
\item the hard violations found by the independent validator (zero is required for publication);
\item room, instructor and group overlaps;
\item the solver outcome (optimal, feasible, infeasible or unknown);
\item wall-clock time.
\end{itemize}
Because parallel search is not deterministic, a different valid timetable can result from the same inputs.

\paragraph{Testing the validator and the gates.} The following checks were used:
\begin{itemize}
\item \textbf{Seeded violations.} Valid timetables were changed to break exactly one rule, and the validator had to report that rule and the sessions it affects.
\item \textbf{Comparison with the solver model.} Given timetables were checked against the solver model by fixing every assignment variable with an equality constraint (not a hint) and solving for feasibility.
\item \textbf{Gate tests.} Automated tests check that an unapproved proposal never reaches the solver input, that an unsupported request creates no rule, that a short time limit on a feasible instance never reports \emph{infeasible}, and that a timetable failing validation cannot be published.
\end{itemize}

\paragraph{Repair diagnosis.} Deliberately impossible inputs were created, such as locked clashes, an instructor unavailable all week, or a daily limit below the teaching load. The diagnosis had to name a minimal set of changes, and a repaired timetable had to pass validation against everything else (Section~\ref{evidence-repair}).

\paragraph{Automated suites.} Backend, frontend and evaluation-harness test suites were run before each reported measurement. Every fault found during evaluation received a regression test.

\section{AI-assisted versus manual entry (RQ3)}\label{meth-rq3}

The proposal planned a within-subject, counterbalanced pilot. Administrators would enter equivalent task lists once with forms and once with typed statements and review. The measures would be completion time, number of corrections, accuracy of the final approved rules against the expected rules, and whether the participant could explain what an approved rule allows. The recording features needed for this exist in the prototype. The study was not run, because no participants or ethics approval were available in the project period. RQ3 is therefore not answered and is listed as future work. No claim about reduced effort is made.

\section{Reproducibility records}\label{reproducibility-records}

Each reported measurement is stored in the repository with the script that produced it, its log and the code commit it ran on (folders under \texttt{thesis/evidence/}, named by date). For the local-model run, every reply is appended to a predictions file as soon as it arrives, together with its latency and token counts. A stopped run therefore continues from the last saved statement. A run note records the commit, the Ollama version, the machine and each restart, and confirms that the interpretation and evaluation code did not change between restarts. The corpus manifest stores hashes of the data, guide, contract and scorer, so the scores can be recomputed from the saved predictions. Fine-tuning runs keep their exported datasets with hashes, the training notebook, the merge check output, every candidate's answers and the rating sheet with its separate key (\texttt{finetuning/}).

\section{Ethics and privacy}\label{ethics-privacy-and-consent}

\begin{itemize}
\item All statements, contexts and training conversations are fictional. The demonstration workspace uses invented names.
\item The language model runs on the same machine as the application. No statement, timetable or personal data was sent to an external AI service during the evaluation.
\item No human participants took part, so no participant data was collected.
\item AI tools were used to draft the corpus, the training examples and parts of the implementation. This is stated wherever such material is evaluated, and the author is responsible for checking the claims and the records.
\end{itemize}

\section{Threats to validity}\label{threats-to-validity}

\paragraph{Construct validity.} The expected answers encode one interpretation and were not reviewed by a second person. Some statements may have more than one reasonable reading. Exact output is a strict measure, so field F1 is reported alongside it. The automatic checks for JAI measure form (language, length, unwanted actions), not truth, and the human rating had a single rater who was not independent.

\paragraph{Internal validity.}
\begin{itemize}
\item Both interpretation methods were adjusted after looking at development errors: the parser's codes and negations, and the model's schema order. Their scores are therefore optimistic.
\item The same author wrote the corpus, the parser, the prompts and the scorer, and also built the solver and the validator. Shared misunderstandings are possible. Seeded faults and the comparison with the solver model reduce this risk but do not remove it.
\item The AI assistant that helped draft the corpus and training data is itself a language model, and its style may favour models of a similar kind.
\end{itemize}

\paragraph{External validity.}
\begin{itemize}
\item The statements are synthetic and may be easier or harder than real administrators' wording.
\item Egyptian Arabic was written by the author and does not cover regional variation.
\item There are two fictional contexts, one machine, one seed and temperature~0.
\item Results concern one 7-billion-parameter model and may not transfer to other models or later versions.
\end{itemize}

\paragraph{Conclusion validity.} Only 36 families and 16--19 assistant cases were available, so the intervals are wide and small differences between methods or rounds should not be over-read.

\section{Summary of limitations}\label{feasibility-risks-and-limitations}

The evaluation is a development-level study: it uses draft labels, synthetic data, a closed set of supported rules and a single locally hosted model family. It shows how the methods behave on this material and whether the surrounding system prevents unreviewed or invalid output from being used. It does not establish accuracy on unseen real statements, a ranking of commercial models, or any reduction in administrative effort.
\clearpage

\clearpage
% ======== sections/implementation.tex ========
\chapter{Implementation}\label{chap:implementation}

% NOTE (migration 23 Nov 2026): Text moved unchanged from main.tex. Much of it was drafted with AI assistance; rewrite it in your own words before submission (the supervisor runs AI checks).
% NOTE (migration 23 Nov 2026): Add an overview of the architecture (frontend, API, worker, database, solver, validator, assistant) with a diagram, in your own words. The status section below dates from 25 October 2026 and is out of date: the requirements appendix has the current status.
\section[Implementation status]{JadwAI implementation status and phased delivery (25 October 2026)}\label{implementation-status}
\emph{Measured evidence from 8 November 2026, including university-scale generation, repair suggestions and the rule-based baseline, is reported in Chapter~\ref{evidence}.}

JadwAI is a new application, separate from the earlier prototype; UniTime is a reference rather than a fork. Its architecture uses React/TypeScript, Python/FastAPI, PostgreSQL, SQLAlchemy/Alembic, a separate scheduling worker, OR-Tools CP-SAT~\cite{google_cpsat} and an independently implemented validator. Phase~4 is complete for the explicitly supported weekly manual scheduling baseline. This milestone does not mean that every proposed SRS variant or the complete research project is finished. Phase~5 evaluation preparation is in progress. JAI currently provides guided interface navigation, not LLM extraction.

\subsection{Academic data, access and approval}
The local foundation includes administrator, viewer and student accounts; password setup/recovery and expiring sessions; login rate limits; restricted database roles; transactional audit records; and optimistic concurrency checks. Normalised academic records cover instructors, rooms, student groups and teaching activities, instructor eligibility, attending-group references and room-equipment requirements. Groups are organised by study year, major and group. Academic CSV imports provide previews, errors, atomic commits and stale-version rejection. Calendar settings are editable; a full parent-course/component hierarchy and rule-replacement history remain unfinished.

The owner-selected calendar has five 90-minute periods: 08:30--10:00, 10:15--11:45, 12:00--13:30, 13:45--15:15 and 15:45--17:15. New runs exclude intervening breaks. Historical drafts retain their original calendars and are not silently relabelled.

Manual rule entry supports mandatory instructor, room and group unavailability; preferred instructor teaching windows; and mandatory activity/group permitted windows. Rules start as drafts and require administrator approval. Approved permitted windows for one target form a union across the week; restrictions on different activity/group targets intersect. Each session must fit entirely inside an allowed window, and a day without a permitted window is excluded for that target. Approved closures prohibit overlapping sessions. The current manual rule schema uses integer-hour boundaries; an exact minute-based natural-language rule must be disclosed as unsupported rather than rounded, even though teaching periods themselves can begin at quarter-hour boundaries.

\subsection{Workload, practical-day policies and assignment locks}
Teaching load is measured in elapsed teaching hours: one 90-minute session is 1.5 hours. It is counted once per instructor, not once per attending group; fixed breaks are excluded. Instructors can have an optional maximum weekly teaching load in half-hour steps. A configured maximum is mandatory; zero prohibits teaching, while a missing value is explicitly unconfigured. The current demonstration leaves weekly maxima unconfigured rather than inventing institutional policy.

Configurable global hard daily caps for students and instructors are implemented, with owner-selected demonstration values of six hours each. The practical-day policy can spread an activity's weekly occurrences across different days, limit consecutive teaching periods and cap internal idle teaching periods. Consecutive periods are adjacent named teaching periods; an internal idle period is an unoccupied teaching period between a resource's first and last session on that day. Fixed breaks are excluded. Preferred daily caps, entity-specific overrides and other proposed institutional variants remain unfinished.

Administrators can copy instructor, day/time and room assignments from a validated saved draft into a draft lock, selecting any nonempty subset. Explicit approval makes the lock mandatory in later generation; it does not modify the source timetable. References are checked against current inputs, contradictory same-target approvals are rejected, and broader contradictions may make generation infeasible. An instructor lock applies to all weekly occurrences of the activity to preserve continuity. Approved locks must be revoked before removal. Input changes mark older drafts stale rather than automatically repairing them.

\subsection{Generation, preliminary diagnostics and independent validation}
Generation runs as a persistent background job. The solver assigns eligible instructors, times and suitable rooms, retaining instructor continuity across each activity's weekly occurrences. Supported hard checks cover completeness and full duration; exact calendar/period fit; capacity, room type and equipment; shared-room, instructor and group non-overlap; approved closures and permitted windows; approved assignment locks; configured weekly maxima; and the configured practical-day policy. Group identifiers represent disjoint cohorts, not individual enrolments. Intervals exclude their end points, so adjoining sessions do not overlap.

On-demand input checks and automatic worker-side preflight distinguish proven local blockers from missing optional limits. Diagnostics identify affected activities, resources and rule identifiers, count available local teaching slots, and detect cases such as unsuitable rooms, exhausted windows, inconsistent locks and excessive group weekly demand. Passing these necessary checks does not establish global feasibility: the solver can still prove the combined resource-allocation problem infeasible.

The independent validator does not import solver model-building code. Saved reports include occurrence/resource/time/rule details for violations, per-check evaluated counts, explicit not-applicable and not-configured states, and the disclosed conflict basis. Structural input failure is labelled not evaluated. An invalid, incomplete or errored candidate cannot become a completed validated draft. This is engineering evidence for the supported vocabulary, not proof of correctness for every conceivable scheduling requirement.

Two versioned objectives are available: minimise sessions outside the assigned instructor's union of preferred windows before minimising internal idle teaching periods, or use a configurable positive-integer weighted sum of these components. A session with applicable windows contributes one preference penalty if it fits none; instructors without such windows contribute no preference penalty. Idle minimisation requires the compact-day policy. The worker independently recalculates preference penalties, idle penalties when active, and the total selected objective before retaining a completed result. Optimality is reported for the selected objective, not for an undefined overall timetable-quality measure.

\subsection{Saved inputs, job outcomes and reproducibility evidence}
Each new run retains an immutable approved-input snapshot with schema version~2, input version, calendar, academic records, approved rules and locks. Canonical input/configuration SHA-256 fingerprints are checked before solving. Saved evidence includes the objective profile and weights, explicit solver parameters, OR-Tools/Python/platform versions, model/application-source fingerprints, validator version, elapsed worker time and queued/started/finished timestamps. These hashes identify recorded content; they do not authenticate a source or capture every hardware and environment detail.

Administrators can set a search bound of 1--60 seconds, a seed and 1--4 search workers. The default is ten seconds, seed~42 and two workers. They may explicitly rerun the selected draft's historical inputs; current changes are not substituted, and the new run remains stale when its input version differs from the workspace. Reusing an idempotency key with a different settings/source payload is refused. Parallel search can produce different feasible schedules despite identical inputs and seed.

Job polling and score-only search progress are implemented. Raw solver outcomes distinguish \texttt{OPTIMAL}, \texttt{FEASIBLE}, \texttt{INFEASIBLE}, \texttt{UNKNOWN} and \texttt{MODEL\_INVALID}; invalid models map to a failed job with an explicit reason. A timeout without a solution is not called proven infeasibility. Cancellation wins when committed before the terminal result write and discards candidates. Interrupted jobs retain an explicit failure reason; the automatic retry limit is zero. Optional retention of a validated cancelled candidate is unavailable. Worker leases and interrupted-job recovery have local tests.

\subsection{Timetable views, student access and exports}
The administrator can display days on the left or across the top and filter by year, major and group. The print layout groups days into rows and teaching periods into columns, accommodates simultaneous sessions and identifies free periods. Schedule PDFs are saved through the browser/system print dialog.

Students can retrieve only their assigned group's timetable. Server-side filtering excludes unrelated activities/groups, manual rules, full administrative snapshots and aggregate statistics; group-permission changes revoke sessions. This is explicitly labelled validated-draft access with stale-input warnings. Publication is unavailable, so this is not the proposed published-only service. Instructor/researcher-specific views and individual-enrolment schedules are not implemented.

Statistics are stored with completed validated runs. Reports are optional downloads, not automatic generation side effects. A one-page PDF uses vector charts for daily sessions, applicable preference satisfaction and resource use. It separately reports room, instructor and group overlap-pair counts; these categories may overlap and are not added as unique conflicts. Unused instructors' preference rules are not applicable rather than satisfied. Room-time use excludes breaks, while occupied-seat use is duration-weighted and excludes unused room periods. These ratios are not an overall efficiency, fairness or educational-benefit score. Reports cover the full saved run; timetable exports respect filters. Suggested filenames describe semester, department, scope/type, draft status, version, date and a short run ID; system print-dialog behaviour still affects final timetable filenames.

\subsection{Phase 5 corpus and offline scorer}
A separate synthetic extraction corpus, distinct from the academic scheduling fixture, is now implemented. Release \texttt{0.1.0-draft} has 128 statements in 32 translated intent families: 96 development and 32 reserved statements, with English, formal Arabic, Egyptian Arabic and mixed language represented equally within each split. The annotation guide and output contract are versioned. Packaging is reproducible; a manifest hashes the data, guide, contract and scorer/authoring code. Development exports omit expected labels. Default operations do not read reserved cases, while explicit reserved access requires a logged reason.

The reserved split is not blind or independently adjudicated. Its author has seen its labels and can access its source. Family-ID and normalized-text checks detect some leakage but do not establish semantic independence. Human multilingual review, supervisor agreement on sampling and definitions, a development pilot, frozen acceptance thresholds and fresh confirmatory data are still required. No model was selected or called, and no paid API run took place.

The offline scorer preserves full-case denominators, flags invalid or missing predictions, compares complete proposal multisets and reports separate decision, clarification, unsupported-content, entity-reference and family-consistency diagnostics. Raw and declared post-review files remain separate. It hashes supplied configuration and predictions and reports only available measured telemetry; absent cost or latency is unknown. Tests that copy expected labels check scorer arithmetic only and must not be reported as AI accuracy findings.

\subsection{Verification evidence and limits}
The latest full application regression run passed 194 backend tests; the frontend suite passed 30 tests and the TypeScript/Vite production build passed. A separate evaluator suite passed 18 tests, including backend schema compatibility for both corpus splits, reproducible packaging, manifest mismatch detection, reserved-access controls, gold-free exports, malformed/duplicate/missing predictions, numeric limits and scoring behaviour. The application and evaluation counts refer to distinct suites; they do not represent independent research observations.

The labelled synthetic scheduling fixture has 64 activities, 28 instructors, 22 rooms, 21 groups and 97 weekly occurrences. The final Phase~4 acceptance run retained all 97 sessions on current input version~12, verified its saved fingerprints and passed independent validation, including recalculation of its selected objective. Its report lists 19 checks, with zero-target checks explicitly not applicable and missing optional weekly maxima disclosed. The live fixture therefore does not demonstrate that every configurable rule was exercised in that one run; seeded and API tests supply additional rule-specific evidence. This is a local engineering example, not a representative benchmark or evidence of superiority over another system.

Additional checks cover small exhaustive scheduling oracles, global infeasibility after passing local checks, invalid-model versus timeout outcomes, concurrency, request-payload retries, cancellation races, failed validation, student scope, report arithmetic and descriptive filenames. The chart PDF was rendered and visually checked, and browser verification exercised input diagnostics, settings and validation coverage. Existing records were backed up before migration~0015. Full backup/restore rehearsal, production container verification, remote CI, deployment, broad cross-browser testing and accessibility conformance remain unestablished.

No extraction benchmark, independent annotation-agreement study or participant study has been conducted. No administrative-effort improvement, scalability claim, learning outcome or real model-performance result is inferred from these tests. Detailed evidence is retained in the project's acceptance, verification and feature-checklist records.

\subsection{Remaining work and phase status}
\begin{enumerate}
\item \textbf{Scope and foundation:} initial design reconciliation, local application foundation and master checklist exist; remote CI, deployment and final institutional acceptance remain pending.
\item \textbf{Academic data and constraint extensions:} core data and the supported manual vocabulary work; parent-course hierarchy, rule-replacement history, preferred/entity-specific daily-limit variants and additional soft types remain unfinished.
\item \textbf{Phase 4 manual scheduling baseline:} complete and locally accepted for its declared supported vocabulary, including workload/daily policies, permitted windows, locks, diagnostics, independent checks and recorded run evidence. This does not complete every broader SRS requirement.
\item \textbf{Phase 5 evaluation preparation:} in progress with the synthetic draft corpus, annotation contract, split/access controls and offline scorer. Human annotation/adjudication, final study sampling, pilot-based thresholds and fresh confirmatory data remain pending.
\item \textbf{Phase 6 JAI extraction:} provider adapter, entity resolution, clarification/review and approval integration remain planned; model selection and spending approval are open.
\item \textbf{Later timetable workflow:} manual scenario editing, comparison, revision-level revalidation, assignment CSV, instructor/researcher views and atomic publication remain unfinished. Publication is absent, not an implemented gate claimed to have passed acceptance.
\item \textbf{Research and thesis completion:} frozen model runs, any approved participant study, analysis, results and discussion remain planned.
\end{enumerate}
Supervisor decisions now concern independent annotation/review, final dataset scale and definitions, development-pilot success thresholds, model/budget selection, participant-study feasibility and the delivery deadline. The teaching-hour unit and six-hour demonstration daily caps have been decided; their suitability as institutional policy still requires agreement. No research result is inferred from implementation or scorer self-tests.


\section[The JAI assistant]{The JAI assistant and its locally hosted model (23 November 2026)}\label{evidence-jai}

\paragraph{What was built.} JAI is a chat panel in JadwAI. Built-in patterns in English, formal Arabic and Egyptian Arabic answer common questions directly. Other messages go to a model served by Ollama on the same machine; nothing is sent to an external provider. The model may explain the supplied facts, discuss trade-offs and prepare at most one website action, which runs only after the administrator confirms it; every proposed action is validated against the same contract as the corresponding page. The base model is \texttt{qwen2.5:7b-instruct} (Qwen2.5-7B-Instruct, 4-bit quantised). On the laptop described above, one answer takes between one and six minutes, so the measurements below concern behaviour, not interactive speed.

\paragraph{Verified guide.} A first evaluation showed that the model invented details about JadwAI (for example a ``third-party provider'' and instructors who ``can edit'' a published timetable). A guide of 25 short statements, each checked against the application code, was therefore added. For each message the (at most three) statements whose keywords match are included in the prompt, and the instructions require the model to use only those facts for questions about JadwAI and otherwise to say that it is not sure. The application also removes repeated sentences, keeps the complete sentences of a reply that was cut off, rewrites a reply once into Egyptian Arabic when an Arabic message received an English reply, and cuts a reply at the first Chinese, Japanese or Korean character (an observed failure of the base model).

\paragraph{Fine-tuning method.} Training examples are fictional conversations in English, formal Arabic and Egyptian Arabic, grouped in \emph{families} (one intent with several wordings) and split into training and held-out families, so that no held-out family is seen in training. The examples were drafted with an AI assistant and approved by the project owner in bulk after reading a change list; this is not an independent item-by-item review. Adapters were trained with QLoRA (rank~16) on Qwen2.5-7B-Instruct on a free cloud T4 GPU, with the loss computed only on JAI's final answer. Because Ollama does not load such adapters for this model, each adapter was merged into Ollama's own 4-bit weights. A first merge back to the original 4-bit format erased about 98\% of the training, because the changes are smaller than the quantisation steps; re-quantising each tensor to a different format (4-bit blocks to 6-bit, 6-bit to 5-bit) preserved it, and a check over every trained matrix found on average 94--97\% of the trained change present per tensor type (lowest single matrix 82\%). Table~\ref{tab:jai-rounds} summarises the training rounds.

\begin{table}[h]
\centering
\small
\begin{tabular}{@{}lp{0.33\textwidth}rrrp{0.22\textwidth}@{}}
\toprule
Round & Training data & Rows & Epochs & Loss & Outcome \\
\midrule
v1 & Original examples, no guide in the prompts & 210 & 2 & 1.62 & Better behaviour, still invents facts \\
v2 & Same answers, guide added to the prompts & 210 & 3 & 1.11 & Worse; not used \\
v3 & Answers rewritten to follow the guide; ``not sure'' examples & 243 & 2 & 1.18 & Preferred in blind rating \\
\bottomrule
\end{tabular}
\caption{JAI fine-tuning rounds (training rows; held-out rows were 69, 69 and 72).}
\label{tab:jai-rounds}
\end{table}


\paragraph{Mixed-direction text.} Arabic replies that contain English names, codes and times were displayed in a confusing order: a range such as 08:30--10:00 inside an Arabic sentence was shown as 10:00--08:30. The chat now chooses each line's direction from its main language and isolates every English or numeric run, which keeps such runs in reading order.
\clearpage

\clearpage
% ======== sections/results.tex ========
% NOTE (migration 23 Nov 2026): Text moved unchanged from main.tex. Much of it was drafted with AI assistance; rewrite it in your own words before submission (the supervisor runs AI checks).
% NOTE (migration 23 Nov 2026): Add the full local-model extraction run (144 statements, thesis/evidence/2026-11-23-full) next to the rule-based baseline when it has finished. Update the test counts in the last section.
\chapter{Results}\label{evidence}

\section{Purpose and boundary}
This chapter records measured evidence about the implemented JadwAI system as of 8 November 2026. It concerns the scheduling engine, its independent validator, the repair diagnosis for impossible timetables, reports and exports, and the interpretation results on the development corpus (RQ1, RQ2). The interpretation results are development-level evidence: the corpus labels remain draft and unreviewed (Section~\ref{dataset-construction}), and no paid provider was used. Engineering checks and research results are kept separate, as in Section~\ref{implementation-status}. Section~\ref{evidence-jai}, added on 23 November 2026, reports the conversational assistant JAI and the fine-tuning of its locally hosted model; those measurements concern assistant behaviour on development cases, not the constraint-extraction questions.

All runs used one laptop (Intel Core i5-1038NG7, 4~cores and 8~threads, 16\,GB memory, macOS~26.5.1), Python~3.14.5 and OR-Tools CP-SAT~9.15.6755~\cite{google_cpsat} with eight search workers and random seed~42. Times are wall-clock seconds on that machine and depend on hardware. Every timetable reported as valid passed JadwAI's independent validator, which recomputes completeness, eligibility, the doctor/TA split, room suitability and capacity, resource overlaps, approved rules, locks, workload and daily policies without using the solver's model. The code is at commits \texttt{e4771a1} and \texttt{4137b96} of the project repository plus the parser correction described in Section~\ref{evidence-baseline}. Scripts and logs are kept in \texttt{thesis/evidence/2026-11-08/}.

\section{Teaching timetable at department scale}\label{evidence-department}
The demonstration workspace contains 448 weekly sessions, 93 student groups in eight programmes, 123 instructors, 65 rooms and 60 advising students (students whose courses differ from their group's, for example a retake, a section change or an elective from another programme). The calendar has six teaching days with five 90-minute periods, a student limit of 7.5 teaching hours per day and at most four consecutive periods.

\begin{itemize}
\tightlist
\item \textbf{One run.} Generation with the convenience objective completed in 253\,s and passed independent validation. None of the 60 advising students has two of their own sessions at the same time. Evaluated on the same timetable without modelling them, 24 of the 60 would have had a clash.
\item \textbf{Published draft.} A later run (222.9\,s) placed all 448 sessions with zero room, instructor and group overlap pairs. All 65 rooms are used by more than one programme, so the absence of room double-booking is a cross-programme property rather than a per-programme one. Room time used was 23.0\% and seats filled 66.2\%, counting advising students who join or drop sessions.
\item \textbf{Two parts.} Built as Applied Computer Science and Informatics (217 sessions) followed by the remaining seven programmes (231 sessions) with the first part fixed, the parts took 128\,s and 129\,s, each passed validation as a combined timetable, and only the complete draft could be published.
\end{itemize}

\section{Teaching timetable at university scale}\label{evidence-university}
The Winter 2026 workbook, entered through the data template, yields 2{,}711 weekly sessions for 349 groups, 669 staff and 65 rooms. This exceeds the 600-session limit of a single run, so the timetable is generated in parts of one or more programme years, each part placed around the fixed sessions of earlier parts (Engineering alone has 912 sessions).

\paragraph{What the input checks found.} Before any search, the input checks identify why the data as entered cannot produce a timetable:
\begin{itemize}
\tightlist
\item 23 of 446 teaching assistants cannot be given an activity of their own, contradicting the rule that every TA teaches;
\item eight lectures for 300--350 students have no room large enough;
\item sessions needing laboratories require 1{,}160 teaching hours a week, while the 18 laboratories provide 810; lecture and tutorial sessions require 2{,}907 hours against 2{,}115;
\item two courses (X118 and ENGN101) each need 87 hours a week from two TAs who have 75 hours between them;
\item every Engineering Year~3 group needs 34.5 of the 36 hours its week can hold under the daily limits (a near-full week rather than an impossibility).
\end{itemize}
The TA names in the entered workbook are placeholders and the 65 rooms appear to cover one building, so these findings describe the entered data rather than the university.

\paragraph{Generation in parts with documented adjustments.} To test the part mechanism at full scale, the experiment set aside the eight oversized lectures and the 23 idle TAs, added 32 placeholder TAs (so no TA pool exceeds half its capacity) and 86 placeholder rooms (25 laboratories, 61 lecture rooms; about 60\% room load), and switched off the daily policies. Table~\ref{tab:parts} reports the result with 120\,s per part.

\begin{table}[h]
\centering
\small
\begin{tabular}{@{}rrrrl@{}}
\toprule
Part & New sessions & Scheduled so far & Time (s) & Validation \\
\midrule
1 & 328 & 328 & 125 & passed \\
2 & 584 & 912 & 127 & passed \\
3 & 572 & 1{,}484 & 127 & passed \\
4 & 581 & 2{,}065 & 128 & passed \\
5 & 598 & 2{,}663 & 130 & passed \\
6 & 40 & 2{,}703 & 128 & passed \\
\bottomrule
\end{tabular}
\caption{University-scale teaching timetable built in six parts (adjusted data; see text).}
\label{tab:parts}
\end{table}

Every earlier session was unchanged in each later part, and the final timetable also passed validation as one whole timetable without the part plan. Two observations qualify the result. First, with the daily policies switched on, Engineering Year~3 alone found no timetable within 300\,s, while without them it was solved in 3\,s; this is consistent with the near-full week reported above. Second, convenience is optimised within each part around fixed sessions, so earlier parts are never improved by later ones; the order of parts matters.

\section{Examination timetable}
For all 203 courses of the workbook, the examination model first computes a proven lower bound on the number of back-to-back examination days that any timetable over the given dates must contain (372). It reached that bound within 20\,s; the previous formulation found no timetable within 20\,s and 423 back-to-back days within 120\,s. Allowing two examinations per day reduced back-to-back days to 302 at the cost of 35 days with two examinations; adding four examination dates reduced them to 51. In the demonstration workspace, 35 examinations over ten days were scheduled in halls with TA invigilators; no advising student has two examinations at the same time, and the one student whose courses cannot be spread further within the dates is named for the administrator.

\section{Repair suggestions for impossible timetables}\label{evidence-repair}
When the approved inputs cannot all hold, JadwAI relaxes what an administrator could change---each approved rule, each fixed session, each instructor's day off and each daily limit---and restores them a chunk at a time while a timetable still exists. The remaining set is minimal (no single member can be restored), not necessarily the smallest possible. The affected sessions are reported from a repaired timetable that the independent validator checks against everything else. Physical data limits (rooms, eligible instructors, ``every TA teaches'', the calendar) are never relaxed. Table~\ref{tab:repair} lists the deliberately impossible cases tested.

\begin{table}[h]
\centering
\small
\begin{tabular}{@{}p{0.42\textwidth}p{0.48\textwidth}@{}}
\toprule
Injected problem & Diagnosis \\
\midrule
Two sessions of one group locked to the same time & Move one session; new day, time and room given; verified \\
The only eligible instructor unavailable all week & Relax that rule; the one affected session listed; verified \\
Instructor's day off is the group's only free day & Either the day off or the group rule (one change); verified \\
Daily limit of 2 hours with 4 hours of teaching in one day & Relax the limit; ``needs 4\,h on Sunday''; verified \\
Two courses on separate days with one teaching day & Relax the rule; both sessions listed; verified \\
Five teaching hours in a four-hour week & No relaxation helps; data must change; overloaded group named \\
\midrule
Demonstration (448 sessions): an instructor unavailable all week and a G7 session locked inside the group's blocked time & Exactly these two causes, the locked session moved to another slot; 49\,s, 19 feasibility solves, 138 sessions examined locally, verified on the whole timetable \\
\bottomrule
\end{tabular}
\caption{Repair diagnosis on impossible timetables.}
\label{tab:repair}
\end{table}

On timetables above 80 sessions a complete diagnosis is too slow (one feasibility solve on the 448-session workspace takes 18--73\,s), so the search starts with the sessions the failure points at and those sharing a group or instructor with them, then proves the fix on the whole timetable within a 240\,s budget. An early version kept a limit that was not needed because a five-second check had ended without a result; inconclusive checks are now retried with more time before a candidate is declared necessary. An inconclusive search is never reported as proof.

\section{Rule-based interpretation baseline}\label{evidence-baseline}
Section~\ref{baselines} describes the two interpretation methods compared for RQ1 and RQ2. JadwAI's built-in parser (deterministic patterns for English, formal Arabic and Egyptian Arabic) was run on all 144 statements of the development release (36 intent families, four language variants each) and scored with the project's offline scorer, which validates each output against the contract and the synthetic context before comparing it with the expected answer. The parser received only entity names as stored in JadwAI (English), not the context's Arabic aliases. A format adapter wrote whole hours as integers (10 rather than 10.0), which the contract requires; half-hour values were left unchanged.

The parser was measured twice. The first measurement exposed that its clarification codes differed from annotation guide~v2 and that it guessed in three situations where the guide requires a question. The parser was then changed, and the second measurement is the one reported in Table~\ref{tab:baseline}. The changes were: (1)~clarification codes renamed to the guide's stable codes (for example \texttt{ambiguous\_instructor} instead of \texttt{ambiguous\_name:ahmed}, and \texttt{unknown\_instructor}); (2)~a mandatory time without a day, a statement naming no rule type (``Nour Hassan, Sunday, 08:00--10:00''), and vague times such as ``early'' now produce questions (\texttt{day\_required}, \texttt{rule\_type\_and\_strength\_required}, \texttt{undefined\_early}) instead of an assumed interpretation; (3)~common Egyptian and English negations that were missing (for example the Egyptian equivalents of ``not available'' and ``cannot'', and ``do not schedule'') were added, together with ``must fit'' and ``no other times''; and (4)~``Monday and Tuesday \ldots{} on each day'' is read as those two days rather than every day. Change~(3) mattered for safety: before it, a statement such as ``Nour Hassan [not available] Sunday 08:00--10:00'' fell through to the ``only'' interpretation and produced the opposite rule. Because these changes were made after inspecting development errors, the reported figures are optimistic for this release and are not evidence of performance on unseen statements.

\begin{table}[h]
\centering
\small
\begin{tabular}{@{}lrrl@{}}
\toprule
Measure & First run & Reported & 95\% interval (reported) \\
\midrule
End-to-end exact output & 0.153 (22/144) & 0.326 (47/144) & 0.215--0.458 \\
Decision accuracy & 0.521 & 0.597 & 0.465--0.736 \\
Proposal precision / recall & 0.222 / 0.315 & 0.448 / 0.361 & \\
Proposal F1 & 0.261 & 0.400 & 0.295--0.515 \\
Field F1 & 0.451 & 0.605 & 0.474--0.712 \\
Clarification F1 & 0.000 & 0.351 & 0.141--0.576 \\
Unsupported-request F1 & 0.000 & 0.000 & \\
\bottomrule
\end{tabular}
\caption{Rule-based baseline on the draft development labels before and after aligning the parser with annotation guide~v2 (intervals: percentile bootstrap over intent families, 1{,}000 replicates).}
\label{tab:baseline}
\end{table}

In the reported run, exact outputs by language were 20/36 (English), 12/36 (mixed), 8/36 (formal Arabic) and 7/36 (Egyptian Arabic); decision accuracy was 0.69, 0.67, 0.53 and 0.50 respectively, and all 144 outputs were contract-valid. Unsupported-request F1 remains zero for two reasons: the parser has no proposals for weekly teaching loads or locks (12 expected), and it deliberately differs from guide~v2 on minute-precise times. The guide labels a closure such as 10:15--11:45 as unsupported; the application stores half hours and rounds closures outwards (10:00--12:00), which suits the 15--30-minute breaks between the institution's teaching periods. This is a recorded design difference, not an error in either. Because the labels are draft and unreviewed, the families are synthetic and the release is development material, these figures describe a baseline on development data only; they are not a confirmatory result and do not compare any language model.

Testing the interpretation workflow on the demonstration workspace also found three faults that were then corrected: a closure ending at 15:15 was stored as ending at 15:00 (closures are now rounded outwards and permitted or preferred windows inwards); the review editor could not show a 10:30 interpretation (it now offers half hours); and Arabic statements could not name instructors stored in English (names are now compared by consonant sounds in both scripts, with Egyptian and classical readings, and a shared first name triggers a clarification).

\section{Reports and exports}
The statistics report for the published draft was recomputed independently from the saved assignments; every figure matched (448/448 sessions, sessions per day, zero overlaps, 65/65 rooms, 123/123 instructors, 22.97\% room time, no extra teaching hours). The review found that seat use ignored advising students, that unset limits were printed as ``None'', and that a report could not be limited to one group; all three were corrected. Timetable PDFs are now generated on the server as landscape A4 pages with short descriptive names (for example ``BA Y3 G7 Schedule'') because browser printing in one environment rotated the page and named the file after the application window.


\section{The JAI assistant: evaluation}\label{evidence-jai-results}

\paragraph{Automatic checks on held-out cases.} Each candidate answered the same held-out cases with the production prompt, at temperature~0. The automatic checks cover only whether the reply is valid, whether it wrongly proposes an action, whether it is in the user's language and whether it stays under 80 words; they do not judge correctness. Table~\ref{tab:jai-auto} reports the conversation cases with the guide in the prompt.

\begin{table}[h]
\centering
\small
\begin{tabular}{@{}lrrrr@{}}
\toprule
16 held-out conversation cases & Base & v1 & v2 & v3 \\
\midrule
Valid reply & 16 & 16 & 16 & 16 \\
No wrongly proposed action & 15 & 16 & 15 & 16 \\
Reply in the user's language & 16 & 13 & 12 & 14 \\
At most 80 words & 14 & 16 & 12 & 14 \\
\midrule
Action choice (5 separate cases) & 3 & 5 & 5 & 4 \\
\bottomrule
\end{tabular}
\caption{Automatic checks. Base, v1 and v2 were run on 22 November; v3 on 23 November with the extended guide, where v1 re-run on identical prompts gave 13 replies in the user's language and 16 under 80 words.}
\label{tab:jai-auto}
\end{table}

Round~2 is a negative result worth recording. Its training answers had been written before the guide existed and were left unchanged when the guide was added to the prompts, so the model was trained to receive the guide and then answer in its own words; it invented more than v1 (``snapshots on the cloud'', solver settings that do not exist). Lower training loss with an extra epoch reflected memorisation of 210 rows rather than better answers. Round~3 rewrote the answers so that every statement about JadwAI comes from the guide entries present in that example's prompt, and the dataset builder now refuses an example whose prompt lacks the guide entries its answer relies on.

\paragraph{Blind human rating.} For v1 and v3, the 19 held-out conversation answers (16 cases and three wordings of a new ``not sure'' question) were placed side by side in random order, without model names, and scored from 1 to 5 on correctness, logic, helpfulness and natural language (Table~\ref{tab:jai-rating}). The rater was the project owner, who had read a short summary of some replies beforehand, so the rating is neither independent nor fully blind.

\begin{table}[h]
\centering
\small
\begin{tabular}{@{}lrr@{}}
\toprule
19 cases, 4 criteria, 1--5 each & v1 & v3 \\
\midrule
Total (maximum 380) & 244 & 274 \\
Mean per case (maximum 20) & 12.8 & 14.4 \\
Correct / logical / helpful / natural (means) & 3.3 / 3.3 / 3.1 / 3.2 & 3.7 / 3.6 / 3.6 / 3.5 \\
Cases won (8 ties) & 3 & 8 \\
\bottomrule
\end{tabular}
\caption{Blind side-by-side rating of fine-tuned v1 and v3.}
\label{tab:jai-rating}
\end{table}

The largest gains were in Arabic: an Arabic privacy question that v1 answered in English, and a request for ``no classes at 08:30'' that v1 answered by claiming such classes were forbidden, while v3 correctly said no mandatory rule prevents them. Six cases received the minimum score for both models, including an off-topic request for a poem and an Arabic question about shared lectures.

\paragraph{Live use.} With v3 in the running application, four questions taken from the weakest cases were asked in the chat. The display of mixed Arabic and English text was correct after a fix described below, but the answers were weak: an Arabic question about laboratories with GPUs used wrong words for rooms; the poem request produced a poem; a frustrated Arabic message drifted into Chinese mid-sentence; and a question about connecting to the university's student information system, asked after the complaint, was answered with a request for a rule identifier from the previous topic. The held-out measurements therefore overstate live reliability, particularly for topic changes within one conversation. These cases were used to write a fourth round of training examples (off-topic requests, topic changes, Arabic room vocabulary and examination publishing), which had not been trained at the time of writing.


\paragraph{Limits.} The examples and contexts are fictional and were drafted with AI assistance; approvals were bulk approvals by the project owner. The held-out sets are small (16--19 conversation cases and 5 action cases), the automatic checks do not measure correctness, and the human rating had one rater who was not independent of the project. All runs used one machine, one random seed and temperature~0. These results describe the development of a local assistant; they are not estimates of constraint-extraction accuracy (RQ1).

\section{Automated tests and limitations of this evidence}
At the time of writing the automated suites pass: 393 backend tests, 71 frontend tests and 47 evaluation-harness tests. They include regression tests for every fault reported in this chapter.

The evidence has the following limits. All measurements come from one machine and one random seed; repeated runs and other hardware may differ, and parallel search can produce different valid timetables from identical inputs. The university-scale run used adjusted data, so it demonstrates the mechanism rather than a timetable for the university. The demonstration workspace is derived from real course structures but uses invented names and advising students. The repair diagnosis returns a minimal rather than a minimum set of changes, depends on time budgets, and on large timetables searches locally first. The interpretation baseline uses draft labels reviewed by no human, a single development release and a parser written by the same author who wrote the corpus. A confirmatory answer to RQ1 and RQ2 would require reviewed labels and a fresh held-out set, and RQ3 would require the participant study outlined in Section~\ref{meth-rq3}.
\clearpage


\clearpage
% ======== sections/conclusion.tex ========
\chapter{Conclusion and Future Work}\label{chap:conclusion}

% TO WRITE (in your own words): answer each research question from the evidence in the Results chapter;
% state the contributions; limitations; future work (participant study if not done, independent annotation,
% round-4 training, larger evaluation). Keep it short: two to four pages.

\section{Summary of findings}
[To be written.]

\section{Contributions}
[To be written.]

\section{Limitations}
[To be written.]

\section{Future work}
[To be written.]
\clearpage

\clearpage
% ======== appendix.tex ========
\appendix
\renewcommand{\appendixtocname}{Appendix}
\renewcommand{\appendixpagename}{\appendixtocname}
\addappheadtotoc
\setboolean{@twoside}{false}
\appendixpage

\chapter{Lists}
\addcontentsline{toc}{section}{List of Abbreviations}
\begin{acronym}[\hspace{3cm}]
  \acro{API}{Application Programming Interface}
  \acro{CP-SAT}{Constraint Programming / Boolean Satisfiability solver (Google OR-Tools)}
  \acro{CSV}{Comma-Separated Values}
  \acro{JAI}{The JadwAI assistant}
  \acro{LLM}{Large Language Model}
  \acro{LoRA}{Low-Rank Adaptation}
  \acro{NFR}{Non-Functional Requirement}
  \acro{FR}{Functional Requirement}
  \acro{QLoRA}{Quantised Low-Rank Adaptation}
  \acro{RQ}{Research Question}
  \acro{SRS}{Software Requirements Specification}
  \acro{TA}{Teaching Assistant}
\end{acronym}
\clearpage
\listoffigures
\addcontentsline{toc}{section}{List of Figures}
\listoftables
\addcontentsline{toc}{section}{List of Tables}

% ---- appendix_content.tex ----
\chapter{Functional and Non-Functional Requirements}\label{app:requirements}

The complete list of requirements, with a one-sentence summary and the prototype status of each, is kept with the source code:

\begin{center}
\url{https://github.com/rawdatarekkkk/jadwai/blob/main/docs/REQUIREMENTS.md}
\end{center}

It is derived from the project's Software Requirements Specification (SRS) v0.4 of 18 November 2026, which defines the inputs, outputs, error behaviour and acceptance criteria of every requirement. SRS v0.2 was reviewed by the supervisor; v0.3 and v0.4 carry the project owner's approval only. The status is the owner's record of the prototype on 23 November 2026, not a supervisor acceptance or a research result: \emph{implemented} means working in the prototype with automated tests, \emph{partial} that a working subset exists.

\section{Functional requirements}

Table~\ref{tab:fr-areas} groups the 78 functional requirements by area. Of these, 55 are implemented, 17 partly implemented and 6 not started.

\begin{longtable}{@{}p{0.26\textwidth}p{0.17\textwidth}p{0.155\textwidth}p{0.33\textwidth}@{}}
\caption{Functional requirements by area.}\label{tab:fr-areas}\\
\toprule
Area & Requirements & Status & Main gaps \\
\midrule
\endfirsthead
\toprule
Area & Requirements & Status & Main gaps \\
\midrule
\endhead
Term data and records & FR-001--FR-008 & Partial & One term only; no aliases or dated availability exceptions \\
Constraint vocabulary and lifecycle & FR-009--FR-013 & Partial & No supersession history or weight editing \\
AI extraction and review & FR-014--FR-023, FR-073, FR-074 & Partial & Extraction quality not yet measured; no audit export with cost \\
Scheduling & FR-024--FR-030, FR-065--FR-070 & Implemented & --- \\
Validation & FR-031--FR-033 & Implemented & --- \\
Scenarios and editing & FR-034--FR-039 & Partial & No explicit-variation scenarios (FR-037) \\
Publication and views & FR-040--FR-043, FR-078 & Implemented & --- \\
Export & FR-044, FR-045 & Implemented & --- \\
Traceability, accounts, audit & FR-046--FR-048 & Partial & Partial requirement-to-publication trace \\
Evaluation support & FR-049, FR-050 & Partial & No human annotations; no participant study \\
Examinations & FR-051, FR-071, FR-072 & Implemented & --- \\
Ask JAI & FR-075--FR-077 & Implemented & --- \\
Infeasibility explanation & FR-058 & Implemented & --- \\
Extensions & FR-052--FR-057 & Mostly not started & Arabic extraction (FR-054) works but is unmeasured \\
Exception handling & FR-059--FR-064 & Partial & Cancelled results are not retained \\
\bottomrule
\end{longtable}

\section{Non-functional requirements}

The SRS defines 23 non-functional requirements and invents no numeric targets; values that depend on later decisions remain open. Table~\ref{tab:nfr} summarises them.

\begin{longtable}{@{}p{0.19\textwidth}p{0.17\textwidth}p{0.56\textwidth}@{}}
\caption{Non-functional requirements and prototype status.}\label{tab:nfr}\\
\toprule
Category & Requirements & Status \\
\midrule
\endfirsthead
\toprule
Category & Requirements & Status \\
\midrule
\endhead
Correctness & NFR-001--003 & Met in tests: separate validator; publication refused on any violation; a time limit is never reported as infeasibility \\
Security & NFR-004, 005 & Server-side permissions; no secrets in the repository. Passwords use scrypt instead of the specified Argon2id (recorded deviation) \\
Privacy & NFR-006, 007 & With the local model no text leaves the server; applicable data-protection rules still to be confirmed \\
Accessibility & NFR-008 & Labelled, keyboard-operable controls; no formal audit \\
Reliability & NFR-009, 010 & Timeouts, retries, malformed output and cancellation tested; recovery objectives open \\
Performance and scale & NFR-011--013 & Runtime, outcome and tested scale recorded with every result; no numeric targets \\
Usability & NFR-014, 015 & Checked manually; not measured (no participant study) \\
Reproducibility & NFR-016 & Versions, seeds and dates recorded for runs and fine-tuning \\
Maintainability & NFR-017 & Separate modules; model provider chosen by configuration \\
Compatibility, retention, backup & NFR-018--020 & Open; local backups before destructive changes \\
Localisation & NFR-021 & English and Arabic interface with right-to-left layout; mixed Arabic/English text kept in reading order \\
Observability, licensing & NFR-022, 023 & Partial; licence list not yet checked \\
\bottomrule
\end{longtable}

\chapter{Screenshots of the Prototype}\label{app:screenshots}

The following screenshots were taken from the running prototype on 23 November 2026 in the demonstration workspace (a browser window of 1440$\times$900 pixels). All names of people are demonstration placeholders. Further screenshots are in the repository folder \texttt{docs/screenshots}.

\begin{figure}[p]
\centering
\thesisimage[width=\textwidth]{02-overview.jpg}
\caption{Overview page of the administrator workspace: counts of activities, instructors, rooms and student groups, the setup path and recent activity.}
\label{fig:shot-overview}
\end{figure}

\begin{figure}[p]
\centering
\thesisimage[width=\textwidth]{03-academic-data.jpg}
\caption{Academic data: teaching activities of each course, with separate tabs for instructors, student groups and rooms, and CSV import.}
\label{fig:shot-academic-data}
\end{figure}

\begin{figure}[p]
\centering
\thesisimage[width=\textwidth]{04-constraints.jpg}
\caption{Constraints: a requirement can be typed in English or Arabic and is interpreted as a proposal; approved rules are listed below with their strength (preference or mandatory).}
\label{fig:shot-constraints}
\end{figure}

\begin{figure}[p]
\centering
\thesisimage[width=\textwidth]{09-proposal-review.jpg}
\caption{Proposal review: a requirement is interpreted by the locally hosted model and shown beside the original text for correction before any rule is created.}
\label{fig:shot-proposal-review}
\end{figure}

\begin{figure}[p]
\centering
\thesisimage[width=\textwidth]{13-teaching-timetable-top.jpg}
\caption{Teaching timetable: generation from approved rules only, building in parts, input checks, search settings and the list of saved drafts.}
\label{fig:shot-teaching-timetable-top}
\end{figure}

\begin{figure}[p]
\centering
\thesisimage[width=\textwidth]{15-teaching-timetable-2.jpg}
\caption{Part of a published teaching timetable: each session shows time, course, activity, instructor and room, with an option to keep it fixed in later regenerations.}
\label{fig:shot-teaching-timetable-2}
\end{figure}

\begin{figure}[p]
\centering
\thesisimage[width=\textwidth]{16-exam-timetable.jpg}
\caption{Examination timetable: midterms by day and time with hall, capacity and invigilator; the summary reports back-to-back examination days across groups.}
\label{fig:shot-exam-timetable}
\end{figure}

\begin{figure}[p]
\centering
\thesisimage[width=\textwidth]{18-jai-chat.jpg}
\caption{The JAI chat panel beside the exam settings: example questions and tasks; a proposed change runs only after the administrator confirms it.}
\label{fig:shot-jai-chat}
\end{figure}

\begin{figure}[p]
\centering
\thesisimage[width=\textwidth]{11-overview-arabic.jpg}
\caption{Arabic interface: the same overview page with right-to-left layout.}
\label{fig:shot-overview-arabic}
\end{figure}

\begin{figure}[p]
\centering
\thesisimage[width=\textwidth]{12-timetable-arabic.jpg}
\caption{Arabic interface: the examination settings page.}
\label{fig:shot-timetable-arabic}
\end{figure}

\clearpage

\chapter{Planning Records from the Proposal}\label{app:planning}

% NOTE (migration 23 Nov 2026): Text moved unchanged from main.tex. Much of it was drafted with AI assistance; rewrite it in your own words before submission (the supervisor runs AI checks).
% NOTE (migration 23 Nov 2026): These sections came from the proposed-methodology chapter. Keep only what still helps the reader, or remove this appendix.
\hypertarget{timeline}{%
\section{Timeline}\label{timeline}}

No deadline has been provided, so no dates are proposed (OQ-14). The
order of work is set out in SRS Section\textasciitilde16 and refined in a
current phased implementation roadmap that places
the annotation guide and evaluation data before AI extraction.

\hypertarget{proposed-deliverables-of-the-study}{%
\section{Proposed deliverables of the
study}\label{proposed-deliverables-of-the-study}}

Constraint vocabulary and schema; annotation guide; annotated dataset
(shareable portion, if permitted); frozen prompts and configurations;
raw outputs; analysis scripts; results with limitations; the prototype;
the thesis.

\hypertarget{consistency-with-the-srs-traceability}{%
\section{Consistency with the SRS
(traceability)}\label{consistency-with-the-srs-traceability}}

\begin{longtable}[]{@{}p{\dimexpr(\linewidth-4\tabcolsep)/2\relax}p{\dimexpr(\linewidth-4\tabcolsep)/2\relax}@{}}
\toprule\noalign{}
Methodology element & Specified feature or research procedure \\
\midrule\noalign{}
\endhead
\bottomrule\noalign{}
\endlastfoot
Extraction with schema-constrained output & FR-015, FR-022 \\
Entity resolution; invented entities & FR-016 \\
Parameters, named periods & FR-001, FR-017 \\
Negation, polarity, hard/soft & FR-018 \\
Clarification and unsupported handling & FR-019 \\
Review; raw versus post-review & FR-020 (stores raw and approved) \\
Approval gate & FR-021 \\
Experiment harness; metrics incl. paraphrase and run-to-run consistency;
versioned development records; frozen held-out configuration; held-out
isolation and access log & FR-049 \\
Annotator rating of clarification-question usefulness; inter-annotator
agreement; sample-size and statistics & research procedures
(dataset/analysis), no FR \\
Audit of model, prompt, cost, latency & FR-023 \\
Manual-entry baseline & FR-010 \\
Manual versus assisted timing and corrections & FR-050 \\
Scheduling measures & FR-024--FR-029, FR-031 \\
Seeded-violation and differential tests of the validator & FR-031 (SRS
EV-7; research procedure) \\
Coverage statements & FR-030, FR-032 \\
System-behaviour tests & SRS AC-01--AC-23 \\
Concurrent changes and publication consistency & FR-059; AC-16 \\
Interrupted operations and recovery & FR-060; AC-17 \\
Duplicate requests and bounded retries & FR-061; AC-18 \\
Incomplete results and cancellation completion & FR-062; AC-19 \\
Validator failures and publication blocking & FR-063; AC-20 \\
Session expiry and permission changes & FR-064; AC-21 \\
Reproducibility records & FR-023--FR-025, FR-046; NFR-016 \\
Differential testing against the solver model & research procedure
(Section\textasciitilde3.11); requires the validator and model to be
separately callable (NFR-002) \\
Rule-based baseline (optional) & research procedure; \textbf{no FR} ---
it is a test-time artefact, not a product feature \\
\end{longtable}

Instructor eligibility, locks and continuity are covered by AC-22; four-language input and the bilingual RTL interface by AC-23. These identifiers refer to the revised scope and do not imply that the older exported SRS has been synchronised. Fixed-assignment solver/validator cross-checks must fix instructor variables too. Infeasibility reporting provides available conflict details; automatic repair remains deferred.

Items marked "research procedure" are deliberately not product
requirements. Every other evaluated capability has an FR.

\hypertarget{decisions-needed-to-finalise-this-chapter}{%
\section{Decisions needed to finalise this
chapter}\label{decisions-needed-to-finalise-this-chapter}}

The remaining decisions concern complete constraint semantics; university template and length; dataset source, size and annotators; success targets; participant study and ethics; model selection and approved cost cap; data/provider policy; institutional workload values and additional daily-limit variants; partial diagnostic result sources; and the submission deadline. Exams, language scope, group-based conflict checking and elapsed-teaching-hour units have been resolved as stated in Chapter~1. The citation style has been resolved as IEEE (BibTeX style \texttt{ieeetr}, numbered in order of first citation). The initial corpus remains unreviewed, so its current size and labels do not resolve the final research sampling or annotation decisions.

\clearpage



% References (IEEE style), generated by BibTeX from overleaf/references.bib.
\renewcommand{\bibname}{References}
\cleardoublepage\phantomsection
\addcontentsline{toc}{chapter}{References}
\begin{thebibliography}{1}

\bibitem{sharma2026schedbench}
S.~S. Sharma and A.~Sharma, ``{SCHEDBench}: A benchmark for evaluating {LLM}
  constraint faithfulness in natural-language combinatorial scheduling.'' arXiv
  preprint arXiv:2608.00991, 2026.
\newblock \url{https://arxiv.org/abs/2608.00991}.

\bibitem{sanguinetti2025preferences}
M.~Sanguinetti, A.~Perniciano, L.~Zedda, A.~Loddo, C.~{Di Ruberto}, and
  M.~Atzori, ``From user preferences to optimization constraints using large
  language models.'' arXiv preprint arXiv:2503.21360, 2025.
\newblock Submitted 27 March 2025; the supplied manuscript carries an
  ITADATA2024 header; final proceedings record not verified.
  \url{https://arxiv.org/abs/2503.21360}.

\bibitem{michailidis2025cpbench}
K.~Michailidis, D.~Tsouros, and T.~Guns, ``{CP-Bench} (later renamed
  {DCP-Bench-Open}): A benchmark of {LLMs} writing constraint models from
  natural-language problem descriptions.'' arXiv preprint arXiv:2506.06052,
  2025.
\newblock Version 2 dated 4 September 2025; title as recorded in the project
  bibliography. \url{https://arxiv.org/abs/2506.06052}.

\bibitem{petrovic2004university}
S.~Petrovic and E.~K. Burke, ``University timetabling,'' in {\em Handbook of
  Scheduling: Algorithms, Models, and Performance Analysis} (J.~Y.-T. Leung,
  ed.), ch.~45, Chapman \& Hall/CRC, 2004.
\newblock Publisher and year not verified from a record.

\bibitem{sobh_course_scheduler}
T.~Sobh, R.~Cousens, and S.~Patel, ``Course {Scheduler}: An automated schedule
  generator.'' University of Bridgeport, n.d.
\newblock Supplied paper; venue and year not verified.

\bibitem{muller2025itc}
T.~M{\"u}ller, H.~Rudov{\'a}, and Z.~M{\"u}llerov{\'a}, ``Real-world university
  course timetabling at the {International Timetabling Competition} 2019,''
  {\em Journal of Scheduling}, vol.~28, pp.~247--267, 2025.
\newblock doi: 10.1007/s10951-023-00801-w.

\bibitem{abdennadher_guc_exams}
S.~Abdennadher and M.~Edward, ``Constraint-based examination timetabling for
  the {German University in Cairo},'' n.d.
\newblock Supplied paper; venue and year not verified.

\bibitem{kaya2026drl}
M.~Kaya, ``Constraint-aware deep {Q}-network with greedy fallback for
  large-scale multi-campus university examination scheduling,'' {\em Array},
  vol.~31, p.~101140, 2026.
\newblock doi: 10.1016/j.array.2026.101140.

\bibitem{google_cpsat}
{Google}, ``{CP-SAT} solver.'' OR-Tools documentation,
  \url{https://developers.google.com/optimization/cp/cp_solver}.
\newblock Accessed 8 October 2026.

\end{thebibliography}

\end{document}
